% !TEX program = xelatex
% 编译方式：xelatex scenario_generation_manual.tex（两遍生成目录）。
\documentclass[UTF8,fontset=fandol,a4paper,11pt]{ctexart}
\usepackage{../../_manual_common/hysim_manual}

\title{\textbf{交直流混合高弹性能源电力系统仿真平台}\\[0.5em]
       {\Large 场景生成技术手册}\\[0.3em]
       {\large 概率模型、随机过程、缩减算法、灾害链与数值验证}}
\author{HySim-XJTU-HRPES（西安交通大学 HRPES 团队）}
\date{\today}

\begin{document}
\maketitle
\begin{abstract}
本手册依据 \srcpath{src/scenario\_generation/}、
\srcpath{include/hacdcpf/analysis/} 与 \srcpath{include/hacdcpf/io/scenario\_bundle\_io.hpp}
的当前实现，系统给出常规、可靠性和弹性三类条件场景的概率口径、确定性基线、二元相关
AR(1) 扰动、气候形变、FMEA 兼容 N--1 目录、台风轨迹/风雨/易损度/故障/修复链、交通影响、
稳健风险特征、尾部锚定混合 k-medoids、JSON/HTTP/工作簿契约及完整验证体系。

本手册的章节组织、理论深度、逐算法源码对应、公开字段覆盖、数值算例和审计边界，明确对标
本项目《潮流计算技术手册》和《最优潮流技术手册》，不再以其他较短模块手册为版式基准。
理论、实现、测试和数值证据分别陈述：独立 Python oracle 只验证公开公式、概率守恒和聚类
结果复算；OpenDSS 与 GridLAB-D 不提供等价的场景生成内核，因而不被虚构为外部认证。
候选概率只在族内归一；气候、易损度、坐标、分段和抢修均有明确适用边界与 fallback 标志。
\end{abstract}

\tableofcontents
\newpage

\section{阅读指南、约定与模块全景}
\label{sec:scenario-overview}

\subsection{手册边界与证据等级}
本手册把内容分为四层，避免把教材理论、代码行为和数值证据混为一谈：
\begin{enumerate}
  \item \textbf{理论层}：给出条件经验分布、随机过程、场景缩减和灾害链的数学定义；
  \item \textbf{源码等价层}：逐项转写当前实现的分支、截断、排序、预算和 fallback；
  \item \textbf{公共契约层}：覆盖公开选项、结果、JSON、HTTP 与工作簿字段；
  \item \textbf{验证层}：区分单元测试、确定性复现、独立公式 oracle 和系统集成测试。
\end{enumerate}
任何“已实现”结论必须同时能指向公开结构、实现函数与测试；任何数值结论必须说明输入、
随机种子、样本数、容差和输出文件。文献中的扩展模型若未进入源码，只能标为理论参考。

\subsection{源码和测试范围}
主要依据如下：
\begin{itemize}
  \item \srcpath{include/hacdcpf/analysis/scenario\_generation.hpp} 与
        \srcpath{src/scenario\_generation/scenario\_generation.cpp}：三族生成、特征与缩减；
  \item \srcpath{include/hacdcpf/analysis/typhoon\_resilience.hpp} 与
        \srcpath{src/scenario\_generation/typhoon\_resilience.cpp}：轨迹、风雨、易损与修复；
  \item \srcpath{include/hacdcpf/analysis/typhoon\_traffic\_impact.hpp} 与同名实现：道路影响；
  \item \srcpath{include/hacdcpf/io/scenario\_bundle\_io.hpp} 与
        \srcpath{src/io/scenario\_bundle\_io.cpp}：场景包验证、规范化和工作簿；
  \item \srcpath{tests/test\_scenario\_generation.cpp}、
        \srcpath{tests/test\_typhoon\_traffic\_impact.cpp}、
        \srcpath{tests/test\_scenario\_bundle\_schema.cpp}：回归合同；
  \item \srcpath{tools/scenario\_generation\_review\_case.cpp} 与
        \srcpath{tests/scenario\_generation\_cross\_validation.py}：固定数值案例和独立复算。
\end{itemize}

\subsection{场景族和入口选择}
\begin{center}\footnotesize
\begin{tabular}{@{}L{2.3cm}L{4.0cm}L{4.2cm}L{4.0cm}@{}}
\toprule
场景族 & 条件信息 & 候选语义 & 公共入口\\
\midrule
Regular & SSP、年份、资产基线、扰动种子 & 多步负荷/新能源/储能 SOC 时序 &
\srcpath{generate\_regular\_scenarios}\\
Reliability & 已实现 FMEA 兼容 N--1 元件目录 & 给定单一元件停运的条件时序 &
\srcpath{generate\_reliability\_scenarios}\\
Resilience & 台风强度、轨迹样本、风雨、故障修复 & 给定灾害强度族的条件事件时序 &
\srcpath{generate\_resilience\_scenarios}\\
组合入口 & 三族 enable 标志 & 分别生成，不推断族间先验 & \srcpath{generate\_scenarios}\\
\bottomrule
\end{tabular}
\end{center}
三个族均先生成候选，再按同一聚类内核缩减；可靠性按 contingency 分组，弹性按 intensity
分组，常规族按 SSP--year 组合分组。组间概率不自动混合。

\subsection{端到端数据流}
\begin{center}
\begin{tikzpicture}[node distance=8mm and 7mm,
box/.style={draw,rounded corners=2pt,align=center,minimum height=8mm,text width=2.7cm},
arr/.style={-{Stealth[length=2mm]}}]
\node[box] (rich) {系统资产\\稳定组件 ID};
\node[box,right=of rich] (base) {资产聚合\\确定性基线};
\node[box,right=of base] (cand) {条件扰动\\候选集合};
\node[box,below=of cand] (risk) {特征、风险源\\尾部/边界锚点};
\node[box,left=of risk] (reduce) {混合 k-medoids\\概率聚合};
\node[box,left=of reduce] (contract) {Result / JSON v3\\HTTP / workbook};
\draw[arr] (rich)--(base); \draw[arr] (base)--(cand);
\draw[arr] (cand)--(risk); \draw[arr] (risk)--(reduce); \draw[arr] (reduce)--(contract);
\end{tikzpicture}
\end{center}
模块本身不执行 PF、OPF、可靠性指标或弹性恢复优化。代表场景进入这些下游模块后，才会
产生电压、潮流、成本、失负荷和恢复时间等物理/决策结果。

\subsection{单位、符号和索引空间}
\begin{itemize}
  \item 功率为 MW/Mvar，能量为 MWh，SOC 和 multiplier 为无量纲；
  \item 时间输入为 h，常规时序步序号不自动代表日历时区；风速 m/s、雨强 mm/h；
  \item 气压亏损 hPa，线路/半径 km，经纬度为度；交通容量保持输入单位；
  \item $N$ 为候选数，$K$ 为代表数，$T$ 为时步数，$S$ 为连续特征维数；
  \item 对外 contingency 与 branch 引用使用稳定组件 \fld{index}；vector position 和图索引
        只可在实现内部使用；AC/DC 同号母线必须保留域限定。
\end{itemize}

\subsection{诚实口径总表}
\begin{center}\small
\begin{tabular}{@{}L{4.3cm}L{5.1cm}L{5.1cm}@{}}
\toprule
对象 & 当前能声明 & 当前不能声明\\
\midrule
概率 & 给定输入和种子下的族内条件经验分布 & 三族联合的无条件发生概率\\
气候 & SSP/year 查表形变及缺失回退 & 区域气候模式或天气预报\\
易损度 & 参数化段级条件概率 & 经现场数据校准的结构可靠度\\
交通 & 边级零维积水与能力递推 & 二维水动力、交通均衡或路由最优性\\
缩减 & 指定距离下的 medoid 局部优化与尾部保留 & 全局最优量化或下游决策误差上界\\
交叉验证 & 独立公式和输出结构复算 & OpenDSS/GridLAB-D 场景生成认证\\
\bottomrule
\end{tabular}
\end{center}

\subsection{复杂度总览}
确定性 profile 和候选生成约为 $O(NT)$；风险特征约为 $O(NS)$；距离矩阵若按需反复求值，
每次完整分配为 $O(NKS)$。当前 medoid 更新枚举簇内候选并计算簇内加权距离，保守最坏上界
为 $O(IKN^2S)$，$I$ 为迭代上限。台风故障链约为 $O(TL_s)$，$L_s$ 为线段数；超过精确抢修
预算时切换到排序近似。空间与候选预算是资源边界，不是统计收敛证书。


% 理论层：对应 PF/OPF 手册的多章理论推导。
\section{场景概率理论与条件经验分布}
\label{sec:scenario-probability-theory}

\begin{theorynote}
场景生成的数学对象不是一组任意曲线，而是给定信息集下、带离散概率质量的有限支持分布。
本章先定义概率口径，再说明风险估计和缩减后概率守恒；实现中的等概率、截断和分组规则见
第~\ref{sec:regular-implementation}--\ref{sec:clustering-implementation} 章。
\end{theorynote}

\subsection{概率空间与信息条件化}
令 $\mathcal I_g$ 表示场景族 $g$ 的输入信息：系统资产、确定性 profile、SSP/year、
contingency 或台风强度、随机种子和算法选项。长度为 $T$ 的场景写为
\begin{equation}
 \omega=(z_0,\ldots,z_{T-1},e),\qquad
 z_t=(P^L_t,P^{PV}_t,P^W_t,P^R_t,E^{st}_t),
\end{equation}
其中 $e$ 可包含故障、修复、轨迹与停运签名。模块输出的是
\begin{equation}
 \widehat P_g(d\omega\mid\mathcal I_g)
 =\sum_{i=1}^{N_g}\pi_i^{(g)}\delta_{\omega_i}(d\omega),
 \quad \pi_i^{(g)}\ge0,
 \quad \sum_i\pi_i^{(g)}=1.
\end{equation}
它是算法产生的条件经验分布，不是现实世界无条件频率。若需要跨族期望，调用方必须提供
$\alpha_g=P(G=g)$，再构造 $P=\sum_g\alpha_g\widehat P_g$；本模块不推断 $\alpha_g$。

\subsection{等概率候选与分层条件概率}
当前三个生成分支在各有效组内采用等概率候选：
\begin{equation}
 \pi_{i\mid h}=\frac1{N_h},
\end{equation}
其中 $h$ 分别代表 SSP--year、contingency 或 typhoon intensity 分组。可靠性结果中
$P(\omega_i\mid C=c)$ 与 $P(C=c)$ 是不同对象；若直接把所有 contingency 下的候选概率
相加，会得到 contingency 数量而不是 1。弹性族同理，不把 TY、STY、SuperTY 的条件样本
误认为类别发生概率。

\subsection{缩减映射与概率守恒}
设 $m_k\in\{1,\dots,N\}$ 为第 $k$ 个 medoid，$A_k$ 是 Voronoi 分配集合：
\begin{equation}
 A_k=\{i:k=\arg\min_j d(\omega_i,\omega_{m_j})\}.
\end{equation}
缩减分布为
\begin{equation}
 \widehat Q_K=\sum_{k=1}^K\Pi_k\delta_{\omega_{m_k}},\qquad
 \Pi_k=\sum_{i\in A_k}\pi_i.
\end{equation}
若 $\{A_k\}$ 构成候选集合的划分，则
\begin{equation}
 \sum_k\Pi_k=\sum_k\sum_{i\in A_k}\pi_i=\sum_i\pi_i=1.
\end{equation}
实现对有限正总质量再做一次归一化，以抵抗浮点累加；数值验证要求误差小于 $10^{-12}$。

\subsection{统计量、风险函数与下游估计}
对任意可积后果 $f(\omega)$，候选与缩减期望分别为
\begin{equation}
 \widehat\mu_N=\sum_i\pi_i f(\omega_i),\qquad
 \widehat\mu_K=\sum_k\Pi_k f(\omega_{m_k}).
\end{equation}
两者差异由距离和 $f$ 对该距离的敏感性共同决定。若 $f$ 对 $d$ 是 $L_f$-Lipschitz，
则有启发式上界
\begin{equation}
 |\widehat\mu_N-\widehat\mu_K|
 \le L_f\sum_k\sum_{i\in A_k}\pi_i d(\omega_i,\omega_{m_k}).
\end{equation}
右侧是 transport mean 的理论动机；但当前代码不估计 $L_f$，因此不能把 transport mean
直接转成 PF/OPF 误差界。

尾部概率和条件尾均值可写为
\begin{align}
 \widehat p_q&=\sum_i\pi_i\mathbf1\{R_i\ge q\},\\
 \widehat{\operatorname{CVaR}}_\alpha(R)
 &=\min_\eta\left[\eta+\frac1{1-\alpha}
   \sum_i\pi_i(R_i-\eta)_+\right].
\end{align}
当前缩减器保留风险分位点和风险源尾部锚点，但不直接优化 CVaR 误差。

\subsection{候选数与蒙特卡洛误差}
对独立同分布样本，均值标准误为 $s/\sqrt N$，Bernoulli 事件概率的近似标准误为
$\sqrt{\widehat p(1-\widehat p)/N}$。但当前候选受条件分组、截断、AR(1) 和确定性 profile
影响，不应机械套用独立样本公式。尤其单条 4096 步 AR(1) 序列的有效样本量近似为
\begin{equation}
 N_{\rm eff}\approx N\frac{1-\rho}{1+\rho},
\end{equation}
当 $\rho=0.75$ 时只有约 $N/7$，这解释了相关系数有限样本偏差为何比白噪声更大。

\subsection{可复现性定义}
固定以下全部条件时，模块要求结果确定性复现：系统资产值及顺序、所有选项、随机种子、
外部气候/SST 资源内容、代码版本和依赖浮点行为。跨编译器或平台若遇精确并列距离，medoid
破同分可能受浮点末位影响；因此可复现性不是跨平台 bitwise 保证。

\begin{gapnote}
当前实现没有为候选分布估计置信区间，也没有自适应增加 $N$ 直到统计收敛；
\fld{candidate\_count} 和目录采样数是预算，不是精度声明。下游若需要概率风险证书，必须另行
执行独立重复、批均值或序贯停止检验。
\end{gapnote}

\subsection{理论--实现对应}
\begin{center}\small
\begin{tabular}{@{}L{5.0cm}L{8.7cm}@{}}
\toprule
理论对象 & 当前实现锚点\\
\midrule
三族条件经验分布 & \srcpath{scenario\_generation.cpp:generate\_regular\_scenarios}、
\srcpath{generate\_reliability\_scenarios}、\srcpath{generate\_resilience\_scenarios}\\
组内等概率 & 三个生成函数中的 candidate probability 赋值\\
medoid 概率聚合 & \srcpath{scenario\_generation.cpp:cluster\_candidates\_core}\\
风险统计与分位点 & \srcpath{scenario\_generation.cpp:enrich\_candidate\_risk\_metadata}\\
跨族先验与置信区间 & \implnone\\
\bottomrule
\end{tabular}
\end{center}

\section{随机过程、相关结构与气候形变理论}
\label{sec:scenario-stochastic-theory}

\begin{theorynote}
本章推导当前二元相关过程及 profile 乘法形变。核心区分是：确定性季节/日内基线描述条件均值，
随机倍率描述围绕该均值的离散不确定性；气候形变改变条件均值，不能与随机噪声混为一项。
\end{theorynote}

\subsection{乘法分解}
对资产族 $a\in\{L,PV,W,R\}$，场景功率写为
\begin{equation}
 P^a_t=P^{a,0}\,s^a_t\,c^a_t(\mathrm{SSP},y)\,m^a_t\,h^a_t,
\end{equation}
其中 $P^{a,0}$ 为资产基准，$s_t$ 为确定性日内/季节形状，$c_t$ 为气候形变，$m_t$ 为随机
倍率，$h_t$ 为台风等事件影响。未启用某层时对应因子取 1。PV 的夜间 $s^{PV}_t=0$，因此
不能通过 $P_t/P^0_t$ 在夜间恢复随机倍率。

\subsection{平稳二元 AR(1)}
令 $x_t=(x^L_t,x^R_t)^\mathsf T$，同一时间相关参数作用于两维：
\begin{equation}
 x_t=\rho_t x_{t-1}+\sqrt{1-\rho_t^2}\,\varepsilon_t,
 \qquad |\rho_t|<1.
\end{equation}
创新由两个独立标准正态变量 $u_t,v_t$ 构造：
\begin{equation}
 \varepsilon^L_t=u_t,\qquad
 \varepsilon^R_t=\rho_{LR}u_t+\sqrt{1-\rho_{LR}^2}\,v_t.
\end{equation}
所以 $\operatorname{Var}(\varepsilon^L)=\operatorname{Var}(\varepsilon^R)=1$，且
$\operatorname{Corr}(\varepsilon^L,\varepsilon^R)=\rho_{LR}$。若初始状态也服从同一二元标准
正态，则
\begin{align}
 \operatorname{Var}(x_t)&=\rho_t^2 I+(1-\rho_t^2)I=I,\\
 \operatorname{Cov}(x_t,x_{t-\ell})&=\rho_t^\ell I,\\
 \operatorname{Corr}(x^L_t,x^R_t)&=\rho_{LR}.
\end{align}
这给出实现所需的目标时间相关和同期交叉相关。

\subsection{块结点与线性插值}
若 \fld{block\_hours} 为 $b>1$，AR(1) 在结点 $j$ 上生成，再对 $t=jb+r$ 做
\begin{equation}
 \widetilde x_t=(1-r/b)x_j+(r/b)x_{j+1}.
\end{equation}
插值改变逐时边际方差和 lag-1 相关；因此 $\rho_t$ 是结点过程参数，不保证插值后每小时样本
相关恰等于输入值。数值相关性验证使用 $b=1$，以隔离 AR(1) 本身。

\subsection{倍率截断与矩偏差}
实现采用仿射映射和硬截断
\begin{align}
 m^L_t&=\operatorname{clip}(1+\sigma_Lx^L_t,l_L,u_L),\\
 m^R_t&=\operatorname{clip}(1+\sigma_Rx^R_t,l_R,u_R).
\end{align}
未触发截断时，$E[m]=1$、$\operatorname{sd}(m)=\sigma$；截断后均值、方差和相关性均发生
偏移。截断概率可由标准正态 CDF 估计：
\begin{equation}
 P_{clip}=\Phi\!\left(\frac{l-1}{\sigma}\right)
 +1-\Phi\!\left(\frac{u-1}{\sigma}\right).
\end{equation}
因此统计测试必须同时报告 multiplier 的最小、最大、均值和样本标准差，不能只检查相关系数。

\subsection{PV 周期性选择偏差}
若仅在 $s^{PV}_t>0$ 时用 $P^{PV}_t/P^{PV,base}_t$ 恢复新能源倍率，观测索引集合
$D=\{t:s^{PV}_t>0\}$ 是确定性的周期性子样本。对自相关过程，
\begin{equation}
 \widehat\rho_D=\operatorname{Corr}(x^L_t,x^R_t\mid t\in D)
\end{equation}
仍以 $\rho_{LR}$ 为总体目标，但有效样本量和有限样本偏差不同。固定算例实测全天风电相关
-0.290252，而 PV 白天子样本为 -0.341279；后者只作诊断，不作为 AR(1) 准入门。

\subsection{储能 SOC 扰动}
对储能设备 $j$，先在 multiplier 空间做最多 128 次截断正态拒绝抽样，再回退 clip：
\begin{equation}
 \widetilde s_j=\operatorname{clip}(s^0_jm_j,s_j^{min},s_j^{max}).
\end{equation}
聚合能量和 SOC 为
\begin{equation}
 E^{st}=\sum_j E_j^{rated}\widetilde s_j,\qquad
 s^{agg}=\frac{E^{st}}{\sum_jE_j^{rated}}.
\end{equation}
当前生成器把该值填充到整个候选时域，不模拟充放电动态；动态 SOC 必须由 time-series
调度/PF 流水线计算。

\subsection{气候形变与随机气候扰动}
月度温升 $\Delta T_m$ 和辐照变化 $\Delta GHI_m$ 来自资源查找；缺失时采用零增量并记录
\fld{climate\_data\_found=false}。对负荷，可概括为
\begin{equation}
 c^L_t=1+\beta_T\Delta T_{m(t)}+\eta^L_{m(t)},
\end{equation}
对 PV 则先改变辐照参考，再施加温度效率项。随机气候扰动使用独立
\fld{climate\_seed}；它与二元 AR(1) 的 seed 空间分离。

\begin{gapnote}
当前模型没有空间高斯场、copula、多变量 VAR、极值尾部拼接、天气状态隐马尔可夫模型，
也没有把 SSP 解释为概率权重。负荷与所有新能源共享一个随机 renewable multiplier，不能
表达 PV、风、其他新能源之间的独立相关矩阵。
\end{gapnote}

\subsection{理论--实现对应}
\begin{center}\small
\begin{tabular}{@{}L{5.0cm}L{8.7cm}@{}}
\toprule
理论条目 & 实现函数与偏差\\
\midrule
二元平稳 AR(1) & \srcpath{scenario\_generation.cpp:correlated\_noise\_factors}\\
块插值与 multiplier clip & 同一函数；插值后逐时相关不等于结点参数\\
确定性 load/PV/wind/other profile & \srcpath{synthesize\_base\_load}、
\srcpath{synthesize\_base\_pv}、\srcpath{synthesize\_base\_wind}、
\srcpath{synthesize\_base\_other\_renewable}\\
储能聚合 SOC & \srcpath{make\_storage\_soc\_profiles}；时域内常值\\
气候查找与形变 & \srcpath{load\_climate\_morph\_factors}、
\srcpath{morph\_regular\_load\_profile}、\srcpath{morph\_regular\_pv}\\
完整多变量天气生成器 & \implnone\\
\bottomrule
\end{tabular}
\end{center}

\section{场景缩减、风险保持与量化理论}
\label{sec:scenario-reduction-theory}

\begin{theorynote}
场景缩减是带离散概率的代表点选择问题。当前实现追求“典型形态与预声明尾部同时保留”，
不是单一运输误差最小化器。本章给出目标、尺度、锚点和评价指标，后续实现章逐分支转写。
\end{theorynote}

\subsection{特征空间}
候选 $i$ 的表示分为三块：连续特征 $a_i\in\mathbb R^S$、停运签名
$\sigma_i\in\{0,1\}^B$ 和 regime 标签 $r_i$。连续特征来自负荷、新能源、净负荷、SOC、
故障概率等摘要，而不是完整 $T$ 维轨迹。该降维降低成本，但可能丢失爬坡、持续时间和
时序相位信息。

\subsection{稳健尺度}
对连续特征列 $s$，令中位数 $m_s$、四分位距 $q_s=Q_{0.75}-Q_{0.25}$，缩放为
\begin{equation}
 \widetilde a_{is}=\frac{a_{is}-m_s}{\max(10^{-9},q_s)}.
\end{equation}
与均值--标准差相比，中位数/IQR 对少量极端值更稳健。当关闭 \fld{robust\_scale} 时使用原始
特征，量纲差异将直接进入距离。

\subsection{混合距离}
\begin{equation}
 d(i,j)=w_c\|\widetilde a_i-\widetilde a_j\|_2
 +w_o\sum_b\mathbf1(\sigma_{ib}\ne\sigma_{jb})
 +w_r\mathbf1(r_i\ne r_j).
\end{equation}
签名长度不同时缺失位按 0；regime 项只在混合方法且双方标签非空时启用。非负权重由 JSON
解析时钳制；距离可能是伪度量，因为不同候选可具有相同摘要、签名和标签。

\subsection{加权 k-medoids 目标}
给定 medoid 集 $M$，基本目标为
\begin{equation}
 J(M)=\sum_i\pi_i\min_{m\in M}d(i,m),\qquad |M|=K.
\end{equation}
分配步使用最近 medoid；更新步在簇成员中枚举新 medoid：
\begin{equation}
 m_k\leftarrow\arg\min_{j\in A_k}\sum_{i\in A_k}\pi_i d(i,j).
\end{equation}
这是交替局部改进，不保证全局最优。冻结 medoid 不参与更新。

\subsection{加权 k-means++ 型初始化}
首个 medoid 和后续候选的选择均受 seed 控制；后续抽样权重为
\begin{equation}
 w_i^{init}=\max(10^{-12},\pi_i)D_i^2,qquad
 D_i=\min_{m\in M}d(i,m).
\end{equation}
它借用 k-means++ 的平方距离思想，但中心必须是候选点，因此仍是 medoid 初始化。

\subsection{风险源构造}
设原始风险源为 $r_{is}$，逐列稳健尺度后得到 $v_{is}$。实现的综合分数为
\begin{equation}
 R_i=\max_s v_{is}
 +0.25\sqrt{\sum_{s:v_{is}>1}(v_{is}-1)^2}
 +0.15\sqrt{S^{-1}\sum_s(v_{is}-\bar v_i)^2}.
\end{equation}
第一项强调单一最坏源，第二项强调多源同时越过一个 IQR，第三项表征源间离散。系数是当前
算法定义，不是概率校准参数。

\subsection{尾部、耦合与决策边界}
对源 $s$，尾部集合为
\begin{equation}
 T_s=\{i:r_{is}\ge Q_{q_s}(r_{\cdot s})\},
\end{equation}
全局尾部对 $R_i$ 同理。若候选至少在两个风险源上超过 coupling quantile，则标记耦合尾部。
决策边界集合为
\begin{equation}
 B_s=\{i:|r_{is}-Q_{q_s}|\le b\max(10^{-9},IQR_s)\},
\end{equation}
其中 $b=$\fld{boundary\_fraction}。它保留阈值附近可能因小扰动改变分类的候选。

\subsection{锚点预算与优先级}
锚点还包括含故障的弹性候选和每个 regime 中风险最高者。若锚点数超过 $K$，实现优先保留
故障候选，再按 $R_i$ 降序截断。这意味着冻结锚点可能占满全部 medoid 预算，导致典型区域
表示变差。\fld{tail\_fraction} 与 \fld{boundary\_fraction} 是不同参数：前者属于尾部配置，
后者是分位点邻域宽度。

\subsection{评价指标}
当前 audit 的主要指标为
\begin{align}
 D_{mean}&=\sum_k\sum_{i\in A_k}\pi_i d(i,m_k),\\
 D_{max}&=\max_k\max_{i\in A_k}d(i,m_k),\\
 C_{global}&=\frac{|T_{global}\cap\{m_k\}|}{|T_{global}|},\\
 C_{source,s}&=\frac{|T_s\cap\{m_k\}|}{|T_s|},\\
 E_{regime}&=\sum_r|p_r-\widetilde p_r|.
\end{align}
$D_{mean}$ 衡量平均运输代理，$D_{max}$ 衡量最坏簇半径，覆盖率衡量尾部代表保留，
$E_{regime}$ 衡量缩减前后标签质量偏移。没有单一指标能证明缩减“更优”。

\subsection{混合方法的不可避免取舍}
冻结尾部锚点对覆盖率单调有利，但可能使 $J(M)$ 变大。本项目 128→12 固定算例中，混合方法
全局尾部覆盖由 0 提高到 1，风电缺额尾部由 0 提高到 1；同时 transport mean 从
1.188380 增至 2.148101，regime mass L1 从 0.046875 增至 1.640625。该结果不是失败，而是
目标冲突的定量证据：若下游关心平均拟合，应选 baseline 或增加 $K$；若关心极端风险，才应
接受更大的平均运输误差。

\subsection{复杂度}
若预先缓存 $N^2$ 距离，内存为 $O(N^2)$，每轮分配 $O(NK)$、medoid 更新最坏 $O(N^2)$；
当前实现按候选计算距离，保守时间上界 $O(IKN^2S)$、额外存储约 $O(NS+NK)$。锚点选择的
分位数排序约 $O(SN\log N)$。

\begin{gapnote}
当前实现没有 Kantorovich/Wasserstein 线性规划、前向/后向概率距离缩减、PAM swap 全邻域、
动态时间规整，也没有用下游 PF/OPF 损失做 decision-aware reduction。审计指标不能替代
下游任务的外样本回测。
\end{gapnote}

\subsection{理论--实现对应}
\begin{center}\small
\begin{tabular}{@{}L{5.1cm}L{8.6cm}@{}}
\toprule
理论条目 & 实现锚点\\
\midrule
特征摘要与稳健尺度 & \srcpath{scenario\_generation.cpp:feature\_summary}、
\srcpath{scaled\_features}\\
混合距离 & \srcpath{scenario\_generation.cpp:candidate\_distance}\\
风险源与综合风险 & \srcpath{raw\_risk\_source\_scores}、
\srcpath{enrich\_candidate\_risk\_metadata}\\
尾部/边界锚点 & \srcpath{select\_tail\_anchors}\\
初始化和 medoid 更新 & \srcpath{initialize\_medoids\_kmeanspp}、
\srcpath{cluster\_candidates\_core}\\
评价指标与基线比较 & \srcpath{compute\_reduction\_metrics}\\
全局最优缩减与下游误差界 & \implnone\\
\bottomrule
\end{tabular}
\end{center}


% 三类场景族与灾害链的源码等价实现。
\section{常规场景：资产聚合、profile 与气候形变}
\label{sec:regular-implementation}

\subsection{入口前置条件与组合展开}
\srcpath{generate\_regular\_scenarios} 对 \fld{candidate\_count}、\fld{cluster\_count} 和
\fld{num\_steps} 取至少 1。SSP 集合优先使用非空 \fld{ssp\_levels}，否则使用单值
\fld{ssp}；年份集合同理。每个 $(ssp,year)$ 组合独立生成候选并聚类，最后把各组合簇拼接到
\fld{RegularScenarioResult.clusters}。因此结果的 \fld{candidate\_count} 是所有组合候选总数，
而不是单组合配置值。

\subsection{负荷资产聚合}
\srcpath{collect\_load\_sites} 的顺序与 fallback 如下：
\begin{enumerate}
  \item 若存在 AC Load，逐个聚合在运负荷 $P=p_{mw}\,scaling$；否则从在运 AC Bus 的
        \fld{pd\_mw} 构造 \texttt{ac\_bus\_load} fallback；
  \item DC 侧同理：优先 DC Load，否则使用 DC Bus 的 \fld{pd\_mw}；
  \item 非正功率站点被丢弃；坐标从对应域母线读取；
  \item 组件 ID 由域和稳定 \fld{index} 构造，不把 AC/DC 同号 bus 合并。
\end{enumerate}
总负荷为
\begin{equation}
 P_L^0=\sum_{j\in\mathcal L_{ac}}\max(0,p_j scaling_j)
       +\sum_{j\in\mathcal L_{dc}}\max(0,p_j scaling_j).
\end{equation}

\subsection{新能源资产分类}
\srcpath{renewable\_breakdown\_mw} 遍历 AC RenewableGen、PVSystem、AC StaticGenerator、
DC PVArray 与 DCStaticGenerator。容量优先使用正的额定/上限字段，否则回退当前出力；静态
发电还乘 \fld{scaling}。按显式 RenewableType 或静态机类型字符串分为 PV、Wind、Other。
\begin{equation}
 P_R^0=P_{PV}^0+P_W^0+P_O^0.
\end{equation}
常规候选分别保存 \texttt{pv\_mw}、\texttt{wind\_mw}、
\texttt{other\_renewable\_mw} 和合计值，避免把所有新能源强制视为 PV。

\subsection{确定性负荷形状}
令 $h=t\bmod24$，$d=\lfloor t/24\rfloor$。源码基线为
\begin{align}
 s_L(h)&=0.84+0.18\sin\frac{2\pi(h-7)}{24}
              +0.10\sin\frac{4\pi(h-18)}{24},\\
 s_{season}(d)&=\begin{cases}
 1+0.10\cos\frac{2\pi(d-213)}{365},&T\ge8760,\\1,&T<8760,
 \end{cases}\\
 P^L_t&=P_L^0\max(0.2,s_L(h))s_{season}(d).
\end{align}
逐负荷 profile 用同一形状乘各资产基线，再按步求和。该曲线是内置参考形状，不是读取到的
AMI 计量序列。

\subsection{PV、风电与其他新能源形状}
PV 基线：
\begin{align}
 g(h)&=\max\left(0,\sin\frac{\pi(h-6)}{12}\right),\\
 s_{PV}(d)&=\begin{cases}0.90+0.18\sin\frac{2\pi(d-80)}{365},&T\ge8760,\\1,&T<8760,
 \end{cases}\\
 P^{PV}_t&=P_{PV}^0 g(h)\max(0.15,s_{PV}(d)).
\end{align}
风电组合日内、天气尺度和季节项，并把低于 0.08 的 availability 置零，否则钳至 1.20；
短于年度时域时 calm-weather 和 season 取 1。其他新能源仅使用弱季节正弦。完整赋值见
\srcpath{synthesize\_base\_pv}、\srcpath{synthesize\_base\_wind} 和
\srcpath{synthesize\_base\_other\_renewable}。

\subsection{气候资源查找}
\srcpath{load\_climate\_morph\_factors} 查找 SSP/year 对应的月度温度和 GHI 增量；成功时 audit
写入 12 个 \fld{climate\_month\_delta\_T}、\fld{climate\_month\_delta\_GHI}、年平均与资源
来源。找不到资源时：
\begin{itemize}
  \item 所有增量置零；
  \item \fld{climate\_data\_found=false}；
  \item source 写为 \texttt{zero\_delta\_missing:ssp:year}；
  \item 若提供 warnings 指针，则返回可见 warning。
\end{itemize}
缺失资源不会被伪装成“已找到零变化气候情景”。

\subsection{负荷分项气候形变}
源码把 regular load profile 分成基础、采暖、制冷等分项并分别形变，再聚合；逐组件 profile
存在时同步修改，只有 aggregate 时保留 fallback。气候随机项由 \fld{climate\_seed} 驱动，
其标准差字段分别控制温度、GHI 和负荷百分比扰动。变形输出及来源写入 candidate audit，
而不是只改变总功率后丢失归因。

\subsection{相关随机倍率的应用顺序}
每个候选以组合索引和候选索引派生 RNG；先构造负荷/新能源相关倍率，再依次作用于各 profile：
\begin{align}
 \widetilde P^L_t&=P^L_t m^L_t,\\
 \widetilde P^{PV}_t&=P^{PV}_t m^R_t,\\
 \widetilde P^W_t&=P^W_t m^R_t,\\
 \widetilde P^O_t&=P^O_t m^R_t.
\end{align}
所有新能源共享 $m^R_t$。若对应 enable 为 false，倍率向量为 1。倍率在 profile 基线为零时
不可由结果相除恢复，这一点对夜间 PV 尤其重要。

\subsection{储能与 standard time series}
在运且额定能量为正的 AC/DC Storage 进入 SOC 聚合。每个候选产生
\texttt{storage\_soc} 和 \texttt{storage\_energy\_mwh}；当前值在时域内保持常数。
\srcpath{build\_standard\_time\_series} 同时生成面向下游的 multiplier profile：负荷、PV、风、
其他新能源和储能均保留 base 值与 capacity 字段。基线近零时使用约定 fallback，并写 warning。

\subsection{候选特征和聚类输入}
\srcpath{make\_candidate\_with\_load\_components} 构造 \fld{TimeSeriesData}、组件负荷 profiles、
standard time series 和 feature summary。特征至少包含 load/PV/wind/renewable/net-load 的能量
或摘要量、SOC 与气候源信息。候选 ID 由 family、SSP、year、序号构成，固定输入和 seed 下稳定。

\subsection{算法流程}
\begin{enumerate}
  \item 聚合负荷、新能源和储能资产；生成确定性 profile；
  \item 展开 SSP/year 组合并读取气候资源；
  \item 对每个候选生成相关 multiplier 和储能 SOC；
  \item 执行气候形变、构造候选、标准 profile 与特征；
  \item 组内赋等概率，计算风险元数据，执行聚类；
  \item 累积组合 audit、warnings、候选数和簇数。
\end{enumerate}

\subsection{预算和大规模 fallback}
年度 $T=8760$ 且候选很多时，完整逐组件 JSON 会迅速膨胀。实现对大规模年度候选使用有界
aggregate profile 路径，audit 明确记录没有丢弃候选但降低了逐组件覆盖。该行为由测试
\texttt{Large annual scenario generation uses bounded aggregate load profiles} 锁定。

\subsection{失败语义与边界}
无资产仍可生成零/基线 profile，但 warning 说明覆盖不足；非法计数被钳到最小 1，而不是抛出。
未知 SSP/year 使用零气候增量。模块不读取实测天气、不拟合负荷预测模型、不执行调度，
也不保证生成 profile 经过 PF 后可行。

\section{可靠性场景：FMEA 兼容 N--1 目录}
\label{sec:reliability-scenarios}

\subsection{条件语义}
可靠性场景描述
\begin{equation}
 \widehat P(d\omega\mid C=c,\mathcal I),
\end{equation}
其中 $C=c$ 是一个已枚举的单元件停运。模块不根据 MTBF/故障率计算 $P(C=c)$；每个
contingency group 内的候选等概率。若下游需要 EENS 或年故障频率，必须从可靠性模块引入
故障率和修复时间。

\subsection{目录来源与资产覆盖}
\srcpath{enumerate\_n1\_contingencies} 复用已实现的可靠性 FMEA 目录和 resolver，而不是单独
发明一套元件枚举。公开类型包含 AC/DC branch、VSC、DC/DC、AC/DC bus、发电机、静态发电、
新能源、PV、负荷、储能、mobile storage、两/三绕组变压器、VPP 与 microgrid。

公开 include 标志控制 AC/DC branch、VSC、DC/DC、bus、generator、load 和 storage 类别。
\fld{include\_vpps} 与 \fld{include\_microgrids} 为兼容字段，默认 false，当前不会把目录扩展到
可靠性 assessment 已实现范围之外。文档不能因为 enum 中存在类型就宣称目录一定产生该类型。

\subsection{稳定 ID 与影响集合}
每个 \fld{ContingencyDefinition} 含：
\begin{itemize}
  \item 稳定字符串 \fld{id}、枚举 \fld{type}、稳定 \fld{component\_index} 和显示名；
  \item 受影响 AC/DC branches、converters、DC/DC converters、buses、generators、loads、storage；
  \item 影响列表使用带域前缀的字符串，避免 AC/DC 同号 ID 冲突。
\end{itemize}
outage signature 根据统一目录位置构造，仅用于聚类 Hamming 距离；对外报告仍返回稳定 ID。

\subsection{目录截断}
若 \fld{max\_contingencies}>0，只保留确定性枚举顺序中的前 $M$ 个目录项。该截断是工作预算，
不是风险排序。大目录时 audit 样本也有界，但完整 contingency coverage 不因 audit 截断而丢失；
分别由测试
\texttt{Large reliability catalog bounds audit samples without dropping contingencies} 与
\texttt{Very large reliability catalogs preserve N-1 coverage within work budgets} 验证。

\subsection{候选构造}
对每个 contingency：
\begin{enumerate}
  \item 使用单步 \fld{num\_steps}=1；JSON 解析完成后也强制为 1；
  \item 生成 \fld{candidates\_per\_contingency} 个负荷/新能源/SOC 扰动候选；
  \item 写入 contingency、outage signature 和 regime label；
  \item 组内赋 $1/N_c$ 概率；
  \item 用最多 \fld{cluster\_count\_per\_contingency} 个 medoid 缩减；
  \item 保存组级 audit，再累计 \fld{cluster\_total}。
\end{enumerate}
可靠性场景生成只标记停运条件，不实际调用 PF、孤岛检测或负荷削减。

\subsection{元件筛选和在运状态}
目录遵循可靠性 resolver 的可停运/在运判定；已停运、无效或不可解析资产不会被静默映射到
另一位置。显示名只用于可读性，身份由类型和稳定 index 决定。

\subsection{概率和报告口径}
若有 $C$ 个 contingency、每组 $N_c$ 个候选，则所有候选 \fld{probability} 的全局和为 $C$，
不是 1：
\begin{equation}
 \sum_{c=1}^C\sum_{i=1}^{N_c}\pi_{i\mid c}=C.
\end{equation}
因此 HTTP/JSON 使用嵌套 group 结构保留条件边界。任何把 groups flatten 后重新归一的调用方
都改变了语义，必须显式提供 contingency prior。

\subsection{复杂度}
设目录大小 $C$，每组候选 $N_c$，特征维数 $S$，代表数 $K_c$。生成约 $O(CN_c)$，聚类
保守最坏 $O(CIK_cN_c^2S)$。\fld{max\_contingencies} 只限制 $C$；它不自动降低每组候选。

\subsection{已知边界}
\begin{gapnote}
当前路径不是 N--k、共因失效、保护拒动、时序 Markov 修复或概率潮流；bus outage 对支路和
设备的拓扑传播依赖 resolver 已实现语义。候选不含故障率权重，不能直接计算年可靠性指标。
\end{gapnote}

\subsection{实现对应}
\begin{center}\small
\begin{tabular}{@{}L{5.2cm}L{8.5cm}@{}}
\toprule
步骤 & 源码锚点\\
\midrule
FMEA 兼容目录 & \srcpath{scenario\_generation.cpp:enumerate\_n1\_contingencies}\\
候选生成与单步强制 & \srcpath{scenario\_generation.cpp:generate\_reliability\_scenarios}\\
JSON options 强制 num steps=1 & \srcpath{scenario\_generation.cpp:scenario\_generation\_options\_from\_json}\\
组级缩减 & \srcpath{scenario\_generation.cpp:cluster\_candidates}\\
故障率和年指标 & \implnone（属于 reliability 模块）\\
\bottomrule
\end{tabular}
\end{center}

\section{弹性场景：强度分组、台风目录与事件定义}
\label{sec:resilience-scenarios}

\subsection{条件语义与分组}
弹性族输出
\begin{equation}
 \widehat P(d\omega\mid H=h,\mathcal I),
\end{equation}
其中 $H$ 是请求的台风强度类别。默认只生成 TY；可选 TD、TS、STS、TY、STY、SuperTY。
每个强度使用 \fld{cluster\_count\_by\_intensity} 的覆盖值，否则用
\fld{default\_cluster\_count}。组间没有类别先验。

\subsection{目录生成}
\fld{TyphoonCatalogOptions} 指定月份范围、每月样本数、基准 seed、时域、步长、随机开关、
月度默认和 SST 资源。\srcpath{build\_typhoon\_catalog} 为每月/样本派生确定性 seed，调用
\srcpath{generate\_typhoon\_track\_sample}，按最大风速分类并建立
\fld{by\_category} 索引。若指定 catalog path，可加载或保存资源；
\fld{loaded\_from\_disk} 和 \fld{source\_path} 保留来源。

\subsection{线程安全缓存}
\srcpath{get\_or\_build\_typhoon\_catalog} 按影响目录内容的选项建立缓存 key，并返回
指向不可变目录的共享指针。混合选项并发刷新必须得到各自一致快照，测试
\texttt{Typhoon catalog snapshots survive concurrent mixed-option refreshes} 防止共享可变目录
被另一线程覆盖。

\subsection{类别采样与 fallback}
\srcpath{sample\_typhoon\_catalog} 先在请求类别索引中按 selection seed 抽样。若类别为空，
实现选择可用的替代类别并置 \fld{used\_category\_fallback=true}；事件同时保留
\fld{requested\_intensity} 和 \fld{selected\_intensity}，不得只返回后者。选中样本的
\fld{selected\_sample\_id}、\fld{selected\_track\_max\_vmax\_ms} 和
\fld{used\_catalog\_sample} 进入结果。

\subsection{轨迹生成与预计算轨迹}
\fld{TyphoonScenarioOptions.use\_precomputed\_track=true} 时直接消费
\fld{precomputed\_track}；否则从初始纬经度、气压亏损、RMW、航向、平移速度、月份和 SST
生成轨迹。预计算路径用于确定性测试和外部轨迹接入，避免为了复算 Holland 核而再次随机生成。

\subsection{事件定义}
\fld{ResilienceEventDefinition} 完整保存：事件 ID、请求/选中强度、样本 ID、fallback 标志、
修复近似标志、故障数、月份、faults、track、branch risks、PV/wind/renewable/load 台风倍率，
以及可选 stage-1 起止步。它是条件场景元数据，不是恢复优化结果。

\subsection{候选生成流程}
对每一 intensity 和候选：
\begin{enumerate}
  \item 派生台风 selection seed，选择目录样本或生成轨迹；
  \item 复用预计算线路分段，计算风雨、支路风险、故障与修复；
  \item 计算 PV、wind、renewable 和 load 台风倍率；
  \item 与常规基线、随机倍率和储能 SOC 合成多步 time series；
  \item 写入事件、outage signature、regime label 与风险源；
  \item 组内赋等概率，执行尾部锚定聚类并记录 audit。
\end{enumerate}

\subsection{设备影响倍率}
PV 风致降额由 \fld{pv\_cutout\_start\_ms}、\fld{pv\_cutout\_full\_ms}、
\fld{pv\_max\_reduction} 和过渡小时控制；风电使用 cut-in/rated/cut-out 和 ramp exponent；
负荷从 \fld{load\_wind\_start} 到 \fld{load\_wind\_full} 线性降至最大削减量。新能源总倍率
按 PV、风、其他容量加权，不把 PV 倍率直接复制给全部新能源。

\subsection{空间预算}
线路自动分段在每个候选前预计算并复用，避免 $N$ 次重复几何构造。每支路段数由
\fld{target\_segment\_length\_km} 推导，并钳在 min/max 之间。没有坐标时使用 case center 与
bus spacing fallback；没有几何长度时产生 synthetic segment。对应标志进入 fault result。

\subsection{概率与强度分类}
最大轨迹风速由 \srcpath{max\_track\_vmax\_ms} 计算；默认可跳过首点以避免初始化值污染。
\srcpath{classify\_typhoon\_intensity} 将最大风速映射到公开类别。分类只是标签，不产生类别
发生概率；目录中的 category count 只描述生成样本数量。

\subsection{复杂度和并行性}
设强度数 $H$、每类候选 $N_h$、轨迹步 $T$、总线段 $L_s$，故障链成本约
$O(\sum_hN_hTL_s)$，随后再加各组聚类成本。目录缓存降低轨迹批量构造成本，但结果仍依赖
选择 seed。线路分段复用是主要空间计算优化。

\subsection{已知边界}
\begin{gapnote}
目录样本不是历史台风重分析集合，也未按登陆率、月份频率或海温后验加权。类别 fallback
只保证可运行，不代表物理相邻类别的概率替代。故障和修复进入场景定义，但不自动反馈到
恢复 MIP、DAE 或交通路由求解。
\end{gapnote}

\subsection{实现对应}
\begin{center}\small
\begin{tabular}{@{}L{5.2cm}L{8.5cm}@{}}
\toprule
步骤 & 源码锚点\\
\midrule
轨迹样本和目录 & \srcpath{typhoon\_resilience.cpp:generate\_typhoon\_track\_sample}、
\srcpath{build\_typhoon\_catalog}\\
目录缓存/分类采样 & \srcpath{get\_or\_build\_typhoon\_catalog}、
\srcpath{sample\_typhoon\_catalog}\\
强度候选与倍率合成 & \srcpath{scenario\_generation.cpp:generate\_resilience\_scenarios}\\
设备台风影响 & \srcpath{scenario\_generation.cpp:wind\_generation\_from\_track}\\
类别先验与历史校准 & \implnone\\
\bottomrule
\end{tabular}
\end{center}

\section{台风轨迹、Holland 风场与降雨核}
\label{sec:typhoon-wind-rain}

\subsection{轨迹点状态}
\fld{TyphoonTrackPoint} 保存 hour、latitude、longitude、Coriolis 参数 $f_c$、中心气压亏损
$\Delta p$、Holland $B$、最大风速半径 $R_{mw}$、航向、平移速度、最大风速和 SST。
这些字段是风雨核的完整公开输入。\fld{vmax\_ms} 用于强度分类；站点风速由 Holland 核重新
计算，二者不可互换。

\subsection{球面距离}
对台风中心 $(\phi_1,\lambda_1)$ 和站点 $(\phi_2,\lambda_2)$，实现使用 Haversine：
\begin{align}
 a&=\sin^2\frac{\phi_2-\phi_1}{2}
 +\cos\phi_1\cos\phi_2\sin^2\frac{\lambda_2-\lambda_1}{2},\\
 r&=2R_E\operatorname{atan2}(\sqrt a,\sqrt{\max(0,1-a)}).
\end{align}
角度先转弧度，$R_E=6371.0088$ km；风雨核将半径下限取 1 km，避免中心奇异。

\subsection{RMW 与 Holland B 参数化}
若 \fld{initial\_rmw\_km}>0，轨迹使用该值；否则
\begin{equation}
 R_{mw}=\exp(3.858-7.7\times10^{-5}\Delta p^2).
\end{equation}
形状参数为
\begin{equation}
 B=\max(0.8,1.881093-0.010917\phi-0.005567R_{mw}).
\end{equation}
求值核另以 $B\ge0.5$、$R_{mw}\ge1$ 防御非法外部轨迹。生成下限与求值下限不同，文档和
测试不得把两者写成同一个值。

\subsection{Holland 梯度风}
记 $q=(R_{mw}/r)^B$、空气密度 $\rho_a=1.15$ kg/m$^3$。源码压力项为
\begin{equation}
 A_p=\frac{B\Delta p\,100}{\rho_a}q e^{-q}.
\end{equation}
令 $r_m=1000r$，$c=f_cr_m/2$，则梯度风为
\begin{equation}
 V_g=\sqrt{\max(0,A_p+c^2)}-c.
\end{equation}
换算至 10 m 阵风：
\begin{equation}
 V_{10}=1.75V_g\left(\frac{10}{275}\right)^{1/7}.
\end{equation}
若 $V_{10}<0.01$ m/s 返回 0，否则返回 $\max(0,V_{10})$。实现没有叠加台风平移速度矢量，
也不输出完整气压廓线。

\subsection{Coriolis 参数}
由纬度生成时
\begin{equation}
 f_c=2\Omega\sin\phi,\qquad \Omega=7.2921\times10^{-5}\ \mathrm{s}^{-1}.
\end{equation}
预计算轨迹可直接提供 \fld{fc}；固定 oracle 使用同一公开输入，不从 C++ 输出反推。

\subsection{降雨径向多项式}
令 $q=\operatorname{clip}(R_{mw}/r,0,1.1)$，基础雨强
\begin{equation}
 r_r=\max(0,-5.5+110q-390q^2+550q^3-250q^4).
\end{equation}
气压强度和变化率修正：
\begin{align}
 k&=\max(1,0.0319\Delta p-0.0395),\\
 \dot p&=\begin{cases}\Delta p_{t-1}-\Delta p_t,&t>0,\\1,&t=0,\end{cases}\\
 k_1&=\max(1,1-\dot p/100).
\end{align}
站点 bearing 与航向相加后归一至 $[0,360)$，按 $45^\circ$ 分成 8 个扇区。扇区倍率同时
根据平移速度（m/s）小于 4、介于 4--8、大于 8 的离散分支选择。最终
\begin{equation}
 R=\max(0,kk_1r_rm_s),
\end{equation}
单位 mm/h。

\subsection{轨迹强度分类}
\srcpath{max\_track\_vmax\_ms} 取轨迹最大 \fld{vmax\_ms}，可跳过首点。分类函数输出
TD、TS、STS、TY、STY、SuperTY 或 Unknown。类别阈值以
\srcpath{typhoon\_resilience.cpp:classify\_typhoon\_intensity} 为准；分类结果只用于目录索引
和分组，不修改站点 Holland 风速。

\subsection{SST 与月度默认}
轨迹生成可使用 SST 资源或月份默认；实际资源路径和月均 SST 进入 track sample/result。
资源不可用时使用内置月度参数化并保留来源。\fld{use\_sst\_resource} 不是“已成功读取 SST”
的证明，必须检查结果 \fld{sst\_resource\_path} 和 sample 元数据。

\subsection{固定数值点}
固定风暴中心 $(22.75,113.55)$、站点 $(22.75,113.95)$、$\Delta p=70$ hPa、$B=1.2$、
$R_{mw}=40$ km、平移速度 18 km/h 时，C++ 返回
\begin{equation}
 V_{10}=55.24244506292464\ \mathrm{m/s},\qquad
 R=32.58499937279819\ \mathrm{mm/h}.
\end{equation}
独立 Python 逐公式复算误差分别为 0 和 $7.82\times10^{-14}$。

\subsection{实现锚点与边界}
\begin{center}\small
\begin{tabular}{@{}L{5.2cm}L{8.5cm}@{}}
\toprule
对象 & 源码锚点\\
\midrule
距离和站点 bearing & \srcpath{typhoon\_resilience.cpp:haversine\_km}、
\srcpath{bearing\_sector\_deg}\\
RMW/B 参数化 & \srcpath{rmw\_from\_delta\_p}、\srcpath{holland\_b}\\
风速 & \srcpath{holland\_wind\_ms\_impl}、\srcpath{holland\_wind\_ms}\\
降雨 & \srcpath{typhoon\_rainfall\_mm\_hr}、\srcpath{rain\_sector\_multiplier}\\
轨迹生成/分类 & \srcpath{build\_track}、\srcpath{classify\_typhoon\_intensity}\\
二维非对称风场、地形/边界层 CFD & \implnone\\
\bottomrule
\end{tabular}
\end{center}

\begin{gapnote}
风雨模型是参数化致灾核，不是业务气象预报。没有地形粗糙度图、台风尺度非对称、降雨雷达
同化或不确定性后验；固定点公式吻合不能推出区域灾害预测精度。
\end{gapnote}

\section{线路分段、易损度、故障抽样与修复}
\label{sec:fragility-fault-repair}

\subsection{自动分段与资产类型}
\srcpath{generate\_typhoon\_line\_segments} 遍历 AC/DC branches。若有母线坐标，沿端点线性插值
构造段；否则围绕 case center 按 bus ID 派生 fallback 坐标，并置
\fld{used\_fallback\_coordinates}。段数近似为支路长度除目标段长，再钳在
\fld{min\_segments\_per\_branch} 与 \fld{max\_segments\_per\_branch}。无法得到实际几何时置
\fld{used\_synthetic\_segments}。

段资产类型为 OverheadLine、Cable 或 Mixed。案例缺少元数据时按
\fld{default\_overhead\_fraction} / \fld{default\_cable\_fraction} 选择默认易损参数；这不是从
线路名称猜测结构类型。

\subsection{风雨等效风速}
架空段先将阵风和降雨组合。令 $v_m=V/1.42$：
\begin{align}
 f_1&=e^{0.006462v_m}-1.2486e^{-0.2769v_m},\\
 f_2&=\begin{cases}0.09376R^{0.7087},&R>0,\\0,&R=0,\end{cases}\\
 V_{eq}&=\max(0,(v_m+\max(0,f_1f_2))1.42).
\end{align}

\subsection{对数正态风易损度}
当 $V_{eq}\le V_d$ 时 $p_w=0$。否则
\begin{align}
 z&=\frac{\ln V_{eq}-\ln V_{50}}{\sigma_{\ln V}},\\
 s&=\begin{cases}
 1,&V_c>0\land V_{eq}\ge V_c,\\
 \Phi(z),&\text{其他},
 \end{cases}\\
 \ell_e&=\operatorname{clip}(\ell,0,5),\\
 p_w&=\operatorname{clip}\left(1-e^{-0.015s\ell_e\Delta t},0,1\right).
\end{align}
$V_c=$\fld{collapse\_wind\_ms} 的饱和分支防止超高风速下数值尾部行为逆转。非正的中位风速或
sigma 由底层 CDF 防御为零 severity。

\subsection{风雨联合段暴露}
令 $r_V=V/V_d$（$V_d\le0$ 时为 0）、$r_R=R/R_d$（$R_d\le0$ 时为 0）：
\begin{align}
 \lambda_s&=\ell^2\exp(a r_V+b r_R+c),\\
 p_s&=\begin{cases}
 \operatorname{clip}(1-e^{-\lambda_s\Delta t},0,1),&V>V_d\lor R>R_d,\\
 0,&\text{其他},
 \end{cases}\\
 p_o&=\operatorname{clip}(1-(1-p_w)(1-p_s),0,1).
\end{align}
该合成假设两个条件机制独立，仅是模型结构，不是由数据估计出的独立性结论。

\subsection{电缆水文与失效概率}
对 Curve Number $CN$：
\begin{align}
 CN'&=\operatorname{clip}(CN,1,99),\quad S=25400/CN'-254,\quad I_a=0.2S,\\
 Q&=\begin{cases}(R-I_a)^2/\max(10^{-9},R-I_a+S),&R>I_a,\\0,&R\le I_a,\end{cases}\\
 I&=\max(0,R-Q),\quad q=Q/(1000\cdot3600),\\
 h&=\left(\frac{q n}{\sqrt{\max(10^{-9},S_0)}}\right)^{0.6},\\
 \theta&=\min\left(\phi,\theta_0+\frac{I}{\max(10^{-9},1000Z)}\right),\\
 p_c&=\operatorname{clip}\left(\frac{p_{c,0}}{10}
 e^{130h+8\theta/\max(10^{-9},\phi)},0,1\right).
\end{align}
Cable 返回 $p_c$，Overhead 返回 $p_o$，Mixed 返回 $1-(1-p_o)(1-p_c)$ 后再钳制。

\subsection{段到支路聚合}
同一支路各段按条件独立并联失效：
\begin{equation}
 p_{b,t}=\operatorname{clip}\left(1-\prod_{s\in b}(1-p_{s,t}),0,1\right).
\end{equation}
结果逐支路保存峰值风速、峰值雨强和峰值失效概率。段越多会提高并联暴露，因此目标段长会
改变支路概率；必须与易损参数共同校准。

\subsection{抽样支持门}
对尚未故障的支路逐步生成 $U\sim U(0,1)$。只有
\begin{equation}
 p_{b,t}\ge p_{min}\quad\text{且}\quad U<p_{b,t}
\end{equation}
时记录首次故障。\fld{min\_fault\_probability} 是抽样支持门，不是风险报告过滤器：低于门槛
的 $p_{b,t}$ 仍进入 branch risk。固定案例把 gate 设为 1，实测峰值概率 0.798910，因此
fault count 为 0，但风险保留。

\subsection{修复时长}
基础 MTTR 优先取资产的正值，否则用 \fld{default\_repair\_time\_hr}：
\begin{align}
 T_r'&=T_r\left[1+k_w\frac{\max(0,V_{peak}-25)}{25}\right],\\
 T_r''&=\operatorname{clip}(T_r',\max(10^{-6},\Delta t),T_{r,max}).
\end{align}
故障事件保存 fault hour 和 repair hour，供下游恢复模块消费。

\subsection{分阶段抢修}
准备时刻为最后故障时刻加非负 delay；维修槽宽
\begin{equation}
 \Delta t_c=\max(\Delta t,\texttt{repair\_crew\_time\_per\_branch\_hr}).
\end{equation}
在精确预算内，每轮试验恢复一条故障支路，依据恢复供电增量及确定性破同分项：
\begin{equation}
 score=impact+10^{-3}p_{peak}+10^{-5}V_{peak}+10^{-6}\ell.
\end{equation}
这是一阶段贪心，不是全时域组合最优。当故障数超过内部精确预算，改按源距离、故障时间、
峰值概率、支路类型和稳定 ID 排序，置 \fld{used\_approximate\_repair\_order=true}。

\subsection{算法流程}
\begin{enumerate}
  \item 生成/接收线路分段；对每个轨迹步计算段中点风雨；
  \item 依资产类型计算段概率并聚合为支路概率；
  \item 更新支路峰值风险，应用 sampling gate 和首次故障抽样；
  \item 计算 MTTR 风速修正；需要时执行 staged repair；
  \item 生成 renewable profile multiplier 和完整 validity flags。
\end{enumerate}

\subsection{复杂度和边界}
风雨/易损链为 $O(TL_s)$。精确 staged repair 的供电增量试验可达 $O(F^2G)$，其中 $F$ 为故障
数、$G$ 为图评估成本；超过预算后排序为 $O(F\log F)$。

\begin{gapnote}
当前没有杆塔级相关失效、飞散物路径、保护动作、维修队空间路由、多队并行资源约束或全局
最优修复调度。分段独立性、默认易损参数和 sampling gate 都必须作为案例假设披露。
\end{gapnote}

\subsection{实现锚点}
\begin{center}\small
\begin{tabular}{@{}L{5.1cm}L{8.6cm}@{}}
\toprule
对象 & 源码锚点\\
\midrule
自动分段 & \srcpath{typhoon\_resilience.cpp:generate\_typhoon\_line\_segments}\\
架空/电缆概率 & \srcpath{overhead\_segment\_probability}、
\srcpath{cable\_segment\_probability}\\
支路聚合与抽样 & \srcpath{generate\_typhoon\_fault\_sequence}\\
修复排序 & \srcpath{apply\_staged\_post\_disaster\_repairs}\\
\bottomrule
\end{tabular}
\end{center}

\section{台风--交通耦合：坐标、积水、速度与容量}
\label{sec:traffic-coupling}

\subsection{输入与输出定位}
\srcpath{apply\_typhoon\_traffic\_impact} 接收 TrafficGraph、台风轨迹和选项，返回受影响图及
逐 link/step 审计。它更新 free-flow time、capacity 和 availability profile，不求解 CTM/LTM、
用户均衡、车辆路径或 MESS 调度。

\subsection{坐标模式}
\begin{itemize}
  \item Geographic：把节点 $x,y$ 解释为 longitude, latitude；
  \item AffineToStudyArea：把网络包围盒仿射映射到指定研究中心和 span；
  \item Auto：若坐标范围合理且网络到轨迹最近距离不超过
        \fld{auto\_geographic\_track\_distance\_km}，使用 Geographic，否则 affine。
\end{itemize}
使用 affine 时置 \fld{used\_affine\_georeferencing=true}。该映射只为暴露计算提供相对空间，
不声称恢复真实道路经纬度。

\subsection{轨迹时间采样}
第 $t$ 个交通步对应 hour
\begin{equation}
 \tau_t=\texttt{start\_hour}+t\Delta t.
\end{equation}
实现选择对应轨迹点并在 link 中点计算 Holland 风和降雨。每条 link 的中点、风、雨均保留在
\fld{TyphoonTrafficLinkStep}，便于独立复算。

\subsection{零维积水递推}
对 link $e$，若存在 per-link drainage/runoff 覆盖则优先使用，否则使用全局值。令
$c_e=c_r m_e$、$d_e$ 为排水率：
\begin{equation}
 W_{e,t+1}=\operatorname{clip}\left(
 W_{e,t}+(c_eR_{e,t}-d_e)\Delta t,0,W_{max}\right).
\end{equation}
单位为 mm。它不含上游汇流、坡面传播或相邻 link 水交换。

\subsection{深度--速度曲线}
源码先计算
\begin{equation}
 v(W)=aW^2+bW+c,qquad
 f_v=\operatorname{clip}\left(\frac{v(W)}{\max(\epsilon,c)},f_{v,min},1\right).
\end{equation}
默认系数来自 Pregnolato 型经验曲线，$W$ 为 mm、$v$ 为 km/h。新 free-flow time 为
\begin{equation}
 \tau_e'=\frac{\tau_e^0}{\max(f_{v,min},f_v)}.
\end{equation}

\subsection{风容量因子}
当 $V\le V_s$，$f_w=1$；否则
\begin{equation}
 f_w=1-\operatorname{clip}\left(
 \frac{V-V_s}{\max(\epsilon,V_f-V_s)},0,1\right)
 (1-\operatorname{clip}(f_{w,min},0,1)).
\end{equation}
综合容量因子为
\begin{equation}
 f_c=\operatorname{clip}(f_v^{\max(0,\gamma)}f_w,f_{c,min},1).
\end{equation}

\subsection{闭锁逻辑}
link 可用当且仅当：原图 link 可用、$W<W_{close}$，且未启用极端风闭锁或 $V<V_{close}$。
闭锁时容量为 0；否则为 $C_e^0f_c$。三个 profile 是否写回分别受
\fld{apply\_free\_flow\_time\_profile}、\fld{apply\_capacity\_profile}、
\fld{apply\_availability\_profile} 控制，计算审计仍保留逐步值。

\subsection{结果聚合}
结果给出至少闭锁一次的 link ID、峰值风/雨/水深、最小速度/容量因子、closed link-step 数和
\fld{model\_scope}。逐步结果的 \fld{link\_index} 是稳定交通 link ID，不是 vector position。

\subsection{固定案例}
固定 4 km、1200 veh/h link，$\Delta t=1$ h，雨强 32.584999 mm/h，runoff 0.85、排水
25 mm/h 时，水深序列为
\begin{equation}
 2.69724947,\ 5.39449893,\ 8.09174840,\ 10.78899787\ \mathrm{mm}.
\end{equation}
对应容量因子为 0.63889991、0.62789772、0.61699343、0.60618705，未达到 300 mm 闭锁。
独立 oracle 最大水深误差 $6.39\times10^{-14}$。

\subsection{输入验证}
实现拒绝非法时间步、负上限、反向 clamp 区间和不物理水文边界；测试
\texttt{Typhoon road impact rejects invalid clamp and hydrology bounds} 锁定这些失败语义。

\subsection{复杂度和边界}
对 $E$ 条 link、$T$ 步，成本 $O(ET)$，输出存储 $O(ET)$。
\begin{gapnote}
当前模型是边级零维水量平衡；不包含二维 Saint--Venant 方程、排水管网、交通需求重分配、
拥堵反馈或抢修车辆可达性自动反馈。affine 坐标只适合缺坐标基准的近似暴露筛查。
\end{gapnote}

\subsection{实现锚点}
\begin{center}\small
\begin{tabular}{@{}L{5.0cm}L{8.7cm}@{}}
\toprule
对象 & 源码锚点\\
\midrule
坐标选择 & \srcpath{typhoon\_traffic\_impact.cpp:georeference\_nodes}\\
风容量因子 & \srcpath{typhoon\_traffic\_impact.cpp:wind\_capacity\_factor}\\
积水/速度/闭锁递推 & \srcpath{typhoon\_traffic\_impact.cpp:apply\_typhoon\_traffic\_impact}\\
CTM/LTM profile 消费 & \srcpath{tests/test\_typhoon\_traffic\_impact.cpp:CTM reads road profiles by parent simulation step}\\
\bottomrule
\end{tabular}
\end{center}

\section{聚类内核：源码等价流程、预算与审计}
\label{sec:clustering-implementation}

\subsection{方法选择}
\fld{method} 等于 \texttt{hybrid\_kmedoids\_tail\_5pct} 或别名
\texttt{hybrid\_kmedoids\_tail} 时进入混合路径；其他字符串使用 weighted k-medoids。
解析器不因未知字符串抛错，因此结果 audit 中的实际 \fld{clustering\_method} 是最终口径。

\subsection{候选数与代表数}
空候选直接返回空簇。实际代表数
\begin{equation}
 K'=\min(\max(1,K),N).
\end{equation}
当 $K'\ge N$ 时每个候选自成一簇，概率保持原值，距离为 0；无需进入迭代。

\subsection{连续特征列顺序}
\srcpath{feature\_summary} 从 profile 计算首值、总量、峰值、缺额和净负荷等候选摘要。
\srcpath{scaled\_features} 按确定性字段名顺序构造矩阵，避免 unordered map 遍历顺序影响距离。
开启 robust scale 时对每列应用 median/IQR；关闭时原值直接进入欧氏距离。

\subsection{风险元数据丰富}
\srcpath{raw\_risk\_source\_scores} 对 regular、reliability、resilience 使用相应风险源：负荷、
净负荷、PV/wind/renewable deficit、故障概率、停运影响等。随后
\srcpath{enrich\_candidate\_risk\_metadata}：
\begin{enumerate}
  \item 对每个源计算分位数、中位数和 IQR；
  \item 写入候选 \fld{risk\_source\_scores} 和综合 \fld{risk\_score}；
  \item 标记 global/source/coupling tail；
  \item 标记 global/source decision boundary；
  \item 汇总 \fld{anchor\_reasons} 并设置 \fld{tail\_anchor}。
\end{enumerate}

\subsection{锚点选择与截断}
\srcpath{select\_tail\_anchors} 合并风险锚点、有故障候选和 regime 极值候选并去重。若超过
$K'$，先按“含故障”布尔降序，再按 risk score 降序，最后按稳定候选顺序破同分，截断到
$K'$。audit 的 \fld{anchor\_ids} 是实际保留锚点，而非截断前全集。

\subsection{冻结语义}
仅当混合方法、\fld{include\_tail\_anchors=true} 且 \fld{freeze\_anchors=true} 时，选中锚点
成为 frozen medoid。冻结中心参与分配但不参与更新。候选本身的 \fld{tail\_anchor} 与
\fld{frozen\_medoid} 是不同字段：前者说明风险原因，后者说明优化状态。

\subsection{初始化}
锚点先进入 medoid 集；剩余位置由加权 k-means++ 型抽样填充。若总抽样权重退化为零，使用
尚未选取的确定性候选 fallback。seed 固定初始化，迭代无随机步骤。

\subsection{分配步}
每个候选计算到所有 medoid 的混合距离，选择严格更小者；完全相等时保留先出现的 medoid。
这使 medoid 排列成为破同分规则的一部分。距离同时使用连续欧氏、outage Hamming 和 regime
惩罚，签名短缺位按 0。

\subsection{更新步}
对每个非冻结非空簇，枚举簇成员 $j$，计算
\begin{equation}
 C_j=\sum_{i\in A_k}\pi_i d(i,j),
\end{equation}
选最小者作为新 medoid。空簇保持原 medoid。若所有 medoid 不变则提前停止，否则最多执行
$\max(1,\texttt{max\_iterations})$ 轮。

\subsection{簇结果}
每簇返回 cluster ID、representative ID、完整代表候选、成员概率和、成员数/ID、加权平均距离、
最大距离、冻结标志和 anchor reasons。平均距离为
\begin{equation}
 \bar d_k=\frac{\sum_{i\in A_k}\pi_i d(i,m_k)}{\Pi_k},
\end{equation}
仅在 $\Pi_k>0$ 时计算。所有簇概率最后按正总和归一。

\subsection{审计指标}
\srcpath{compute\_reduction\_metrics} 计算 transport mean/max、全局/源/耦合尾部覆盖和 regime
mass L1。\fld{compare\_baseline=true} 时，对同一候选再运行无锚点 weighted k-medoids，并保存
两组指标及 delta。它是同输入同距离的算法消融，不是与外部软件对照。

\subsection{固定 128→12 复算}
固定 seed=17、\fld{regime\_weight}=0、\fld{boundary\_fraction}=0.10 时，混合方法 12 个 medoid
全部冻结。独立 Python 从序列化候选重新计算 robust scale、成员距离、加权平均半径和全局运输
指标，最大误差为 0（容差 $10^{-10}$）。结果揭示：
\begin{center}\small
\begin{tabular}{@{}L{4.4cm}rr@{}}
\toprule
指标 & Hybrid tail-anchor & Weighted baseline\\
\midrule
transport mean & 2.148101 & 1.188380\\
transport max & 7.169293 & 3.061194\\
global tail coverage & 1.000000 & 0.000000\\
coupling tail coverage & 0.450000 & 0.150000\\
wind deficit tail coverage & 1.000000 & 0.000000\\
regime mass L1 & 1.640625 & 0.046875\\
\bottomrule
\end{tabular}
\end{center}

\subsection{选择建议}
\begin{itemize}
  \item 平均 profile 重构优先：weighted baseline 或提高 $K$；
  \item 极端后果筛查优先：hybrid，并同时报告 transport/regime 退化；
  \item 停运组合重要：提高 outage Hamming weight；
  \item regime 质量重要：提高 regime weight，但必须检查 transport 变化；
  \item 特征量纲不一致：开启 robust scale；已有物理归一化时可关闭以做审计复算。
\end{itemize}

\subsection{失败和边界}
聚类不返回“收敛到全局最优”的状态；达到迭代上限只表示停止。结果没有下游 PF/OPF 后果误差。
\begin{gapnote}
锚点数量等于 $K$ 时，所有中心被冻结，算法事实上不再优化平均距离。必须通过
\fld{frozen\_medoid\_count}、transport 和 baseline comparison 识别该边界。
\end{gapnote}


% 公共契约、验证与准入。
\section{公共选项与结果结构完整参考}
\label{sec:scenario-options-results}

本章以公开头文件为准覆盖全部字段。默认值来自结构体初始化；JSON 解析会对部分字段钳制或
强制覆盖，另行注明。布尔、枚举、ID 和计数无量纲；典型范围不是自动校验承诺。

\subsection{RegularScenarioOptions}
\compmeta{RegularScenarioOptions}{include/hacdcpf/analysis/scenario\_generation.hpp}
\begin{paramtable}
\fld{enabled} & -- & -- & bool & true & 是否生成常规族\\
\fld{candidate\_count} & $N$ & 个/组合 & $\ge1$ & 100 & 每个 SSP--year 组合候选数；实现至少取 1\\
\fld{cluster\_count} & $K$ & 个/组合 & $\ge1$ & 10 & 每组合代表数；不超过候选数\\
\fld{num\_steps} & $T$ & 步 & $\ge1$ & 8760 & 时序长度\\
\fld{ssp} & -- & -- & 字符串 & ssp245 & 单一 SSP fallback\\
\fld{year} & $y$ & 年 & 整数 & 2050 & 单一年份 fallback\\
\fld{ssp\_levels} & -- & -- & 字符串数组 & 空 & 非空时覆盖单一 SSP\\
\fld{years} & -- & 年数组 & 整数数组 & 空 & 非空时覆盖单一年份\\
\end{paramtable}

\subsection{ReliabilityScenarioOptions}
\compmeta{ReliabilityScenarioOptions}{include/hacdcpf/analysis/scenario\_generation.hpp}
\begin{paramtable}
\fld{enabled} & -- & -- & bool & true & 是否生成可靠性族\\
\fld{candidates\_per\_contingency} & $N_c$ & 个 & $\ge1$ & 30 & 每个 N--1 条件候选数\\
\fld{cluster\_count\_per\_contingency} & $K_c$ & 个 & $\ge1$ & 5 & 每故障代表数\\
\fld{num\_steps} & -- & 步 & 固定 & 1 & JSON 解析后强制为 1\\
\fld{include\_ac\_branches} & -- & -- & bool & true & 包含 AC 支路\\
\fld{include\_dc\_branches} & -- & -- & bool & true & 包含 DC 支路\\
\fld{include\_vsc\_converters} & -- & -- & bool & true & 包含 VSC\\
\fld{include\_dcdc\_converters} & -- & -- & bool & true & 包含 DC/DC\\
\fld{include\_ac\_buses} & -- & -- & bool & true & 包含 AC 母线\\
\fld{include\_dc\_buses} & -- & -- & bool & true & 包含 DC 母线\\
\fld{include\_generators} & -- & -- & bool & true & 包含已实现发电类别\\
\fld{include\_loads} & -- & -- & bool & true & 包含 AC/DC 负荷\\
\fld{include\_storage} & -- & -- & bool & true & 包含已实现储能类别\\
\fld{include\_vpps} & -- & -- & bool & false & API 兼容字段；不扩展 FMEA 目录\\
\fld{include\_microgrids} & -- & -- & bool & false & API 兼容字段；不扩展 FMEA 目录\\
\fld{max\_contingencies} & $C_{max}$ & 个 & $\ge0$ & 0 & 0 表示不截断；正值按枚举顺序截断\\
\end{paramtable}

\subsection{ResilienceScenarioOptions}
\compmeta{ResilienceScenarioOptions}{include/hacdcpf/analysis/scenario\_generation.hpp}
\begin{paramtable}
\fld{enabled} & -- & -- & bool & true & 是否生成弹性族\\
\fld{intensity\_levels} & -- & -- & enum 数组 & TY & 请求强度列表\\
\fld{candidates\_per\_intensity} & $N_h$ & 个 & $\ge1$ & 30 & 每强度候选数\\
\fld{default\_cluster\_count} & $K_h$ & 个 & $\ge1$ & 5 & 无类别覆盖时代表数\\
\fld{cluster\_count\_by\_intensity} & -- & -- & map & 空 & 类别到代表数覆盖\\
\fld{num\_steps} & $T_h$ & 步 & $\ge1$ & 48 & 灾害场景时域\\
\fld{month} & $m$ & 月 & $1\sim12$ & 8 & 轨迹目录月份\\
\end{paramtable}

\subsection{PerturbationOptions}
\compmeta{PerturbationOptions}{include/hacdcpf/analysis/scenario\_generation.hpp}
\begin{paramtable}
\fld{enable\_load\_perturbation} & -- & -- & bool & true & 启用负荷随机倍率\\
\fld{enable\_renewable\_perturbation} & -- & -- & bool & true & 启用共享新能源随机倍率\\
\fld{seed} & -- & -- & uint & 1 & AR(1)/SOC 候选种子基值\\
\fld{load\_sigma} & $\sigma_L$ & pu & $\ge0$ & 0.08 & 负荷倍率未截断标准差\\
\fld{renewable\_sigma} & $\sigma_R$ & pu & $\ge0$ & 0.12 & 新能源倍率未截断标准差\\
\fld{load\_min\_multiplier} & $l_L$ & pu & $\ge0$ & 0.75 & 负荷倍率下限\\
\fld{load\_max\_multiplier} & $u_L$ & pu & $\ge l_L$ & 1.25 & 负荷倍率上限\\
\fld{renewable\_min\_multiplier} & $l_R$ & pu & $\ge0$ & 0.0 & 新能源倍率下限\\
\fld{renewable\_max\_multiplier} & $u_R$ & pu & $\ge l_R$ & 1.20 & 新能源倍率上限\\
\fld{enable\_storage\_soc\_perturbation} & -- & -- & bool & true & 启用储能初始 SOC 扰动\\
\fld{storage\_soc\_sigma} & $\sigma_s$ & pu & $\ge0$ & 0.10 & JSON 解析钳至非负\\
\fld{storage\_soc\_min\_multiplier} & $l_s$ & pu & $\ge0$ & 0.75 & JSON 解析钳至非负\\
\fld{storage\_soc\_max\_multiplier} & $u_s$ & pu & $\ge0$ & 1.25 & 若小于下限则交换\\
\fld{enable\_climate\_perturbation} & -- & -- & bool & true & 启用气候随机扰动\\
\fld{climate\_seed} & -- & -- & uint & 2030 & 气候扰动独立种子\\
\fld{climate\_temp\_sigma\_c} & -- & $^\circ$C & $\ge0$ & 0.7 & 温度扰动标准差\\
\fld{climate\_ghi\_sigma\_pct} & -- & \% & $\ge0$ & 8.0 & GHI 扰动标准差\\
\fld{climate\_load\_sigma\_pct} & -- & \% & $\ge0$ & 3.0 & 气候负荷扰动标准差\\
\fld{block\_hours} & $b$ & h & $\ge1$ & 24 & AR 结点间隔；实现至少取 1\\
\fld{temporal\_correlation} & $\rho_t$ & -- & $[-0.999,0.999]$ & 0.75 & JSON 解析钳制\\
\fld{load\_renewable\_correlation} & $\rho_{LR}$ & -- & $[-0.999,0.999]$ & -0.25 & JSON 解析钳制\\
\end{paramtable}

\subsection{TyphoonImpactOptions}
\compmeta{TyphoonImpactOptions}{include/hacdcpf/analysis/scenario\_generation.hpp}
\begin{paramtable}
\fld{pv\_transition\_hours} & -- & h & $\ge0$ & 3 & PV 降额平滑过渡\\
\fld{load\_wind\_start} & $V_{Ls}$ & m/s & 非负 & 25.0 & 负荷降额起点\\
\fld{load\_wind\_full} & $V_{Lf}$ & m/s & $>V_{Ls}$ & 50.0 & 达最大负荷降额风速\\
\fld{load\_max\_reduction} & $r_L$ & pu & $0\sim1$ & 0.40 & 最大负荷削减比例\\
\fld{wind\_cut\_in\_ms} & $V_{ci}$ & m/s & 非负 & 2.5 & 风机切入\\
\fld{wind\_rated\_ms} & $V_r$ & m/s & $>V_{ci}$ & 11.2 & 风机额定风速\\
\fld{wind\_cut\_out\_ms} & $V_{co}$ & m/s & $>V_r$ & 25.0 & 风机切出\\
\fld{wind\_ramp\_exponent} & $\gamma_w$ & -- & 正数 & 2.0 & 切入到额定功率曲线指数\\
\end{paramtable}

\subsection{ClusteringOptions}
\compmeta{ClusteringOptions}{include/hacdcpf/analysis/scenario\_generation.hpp}
\begin{paramtable}
\fld{method} & -- & -- & 字符串 & {\scriptsize\shortstack[l]{hybrid\_\\kmedoids\_\\tail\_5pct}} & 方法选择\\
\fld{max\_iterations} & $I$ & 轮 & $\ge1$ & 50 & 实现至少执行一轮\\
\fld{seed} & -- & -- & uint & 1 & medoid 初始化种子\\
\fld{continuous\_weight} & $w_c$ & -- & $\ge0$ & 1.0 & JSON 钳至非负\\
\fld{outage\_hamming\_weight} & $w_o$ & -- & $\ge0$ & 1.0 & 停运签名权重\\
\fld{robust\_scale} & -- & -- & bool & true & median/IQR 尺度\\
\fld{include\_tail\_anchors} & -- & -- & bool & true & 选择尾部/边界锚点\\
\fld{regime\_weight} & $w_r$ & -- & $\ge0$ & 1.0 & regime 不同惩罚\\
\fld{tail\_fraction} & -- & -- & $[0,0.5]$ & 0.05 & JSON 钳制；与 boundary fraction 不同\\
\fld{source\_tail\_quantile} & $q_s$ & -- & $[0.5,0.999]$ & 0.95 & 源/全局尾阈值\\
\fld{coupling\_quantile} & $q_c$ & -- & $[0.5,0.999]$ & 0.90 & 多源耦合阈值\\
\fld{boundary\_fraction} & $b$ & IQR & $[0,0.5]$ & 0.02 & 分位点邻域宽度\\
\fld{freeze\_anchors} & -- & -- & bool & true & 锚点是否冻结\\
\fld{compare\_baseline} & -- & -- & bool & true & 同候选运行 baseline 消融\\
\end{paramtable}

\subsection{TyphoonScenarioOptions：轨迹与空间}
\compmeta{TyphoonScenarioOptions}{include/hacdcpf/analysis/typhoon\_resilience.hpp}
\begin{paramtable}
\fld{horizon\_hours} & $H$ & h & $\ge1$ & 48 & 轨迹时域\\
\fld{time\_step\_hr} & $\Delta t$ & h & $>0$ & 1.0 & 时间步\\
\fld{seed} & -- & -- & uint & 2030 & 轨迹/故障抽样种子\\
\fld{stochastic} & -- & -- & bool & false & 是否随机轨迹\\
\fld{use\_precomputed\_track} & -- & -- & bool & false & 使用外部轨迹向量\\
\fld{precomputed\_track} & -- & -- & point 数组 & 空 & 预计算轨迹\\
\fld{requested\_category} & -- & -- & enum & Unknown & 请求强度\\
\fld{selected\_category} & -- & -- & enum & Unknown & 实际选择强度\\
\fld{selected\_track\_max\_vmax\_ms} & -- & m/s & 非负 & 0 & 选中轨迹最大 vmax\\
\fld{selected\_sample\_id} & -- & -- & 字符串 & 空 & 目录样本 ID\\
\fld{used\_category\_fallback} & -- & -- & bool & false & 类别是否回退\\
\fld{month} & $m$ & 月 & $1\sim12$ & 8 & 月份\\
\fld{use\_month\_defaults} & -- & -- & bool & false & 使用月度默认参数\\
\fld{use\_sst\_resource} & -- & -- & bool & true & 尝试使用 SST 资源\\
\fld{sst\_resource\_path} & -- & -- & 路径 & 空 & SST 资源路径\\
\fld{initial\_latitude} & $\phi_0$ & deg & $[-90,90]$ & 20.5 & 初始纬度\\
\fld{initial\_longitude} & $\lambda_0$ & deg & $[-180,180]$ & 114.5 & 初始经度\\
\fld{initial\_delta\_p\_hpa} & $\Delta p_0$ & hPa & 正数 & 44.0 & 初始气压亏损\\
\fld{initial\_rmw\_km} & $R_0$ & km & $\ge0$ & 0 & 0 时由气压推导\\
\fld{initial\_heading\_deg} & -- & deg & 任意 & 0 & 初始航向\\
\fld{initial\_translation\_speed\_kmph} & -- & km/h & 非负 & 21.6 & 平移速度\\
\fld{auto\_segment\_lines} & -- & -- & bool & true & 自动线路分段\\
\fld{target\_segment\_length\_km} & $\ell_s$ & km & 正数 & 0.5 & 目标段长\\
\fld{min\_segments\_per\_branch} & -- & 个 & $\ge1$ & 1 & 每支路最小段数\\
\fld{max\_segments\_per\_branch} & -- & 个 & $\ge1$ & 40 & 每支路最大段数\\
\fld{fallback\_case\_center\_latitude} & -- & deg & 合法纬度 & 22.75 & 缺坐标案例中心\\
\fld{fallback\_case\_center\_longitude} & -- & deg & 合法经度 & 113.55 & 缺坐标案例中心\\
\fld{fallback\_bus\_spacing\_km} & -- & km & 正数 & 0.8 & 合成坐标间距\\
\fld{default\_overhead\_fraction} & -- & pu & $0\sim1$ & 1.0 & 缺资产元数据默认比例\\
\fld{default\_cable\_fraction} & -- & pu & $0\sim1$ & 0.0 & 缺资产元数据默认比例\\
\fld{default\_overhead\_fragility} & -- & -- & struct & 见下 & 架空默认易损参数\\
\fld{default\_cable\_fragility} & -- & -- & struct & 见下 & 电缆默认易损参数\\
\end{paramtable}

\subsection{TyphoonScenarioOptions：故障、修复与新能源}
\begin{paramtable}
\fld{min\_fault\_probability} & $p_{min}$ & pu & $0\sim1$ & 0.05 & 抽样支持门；风险仍报告\\
\fld{default\_repair\_time\_hr} & $T_r$ & h & 正数 & 6.0 & 资产无 MTTR 时 fallback\\
\fld{repair\_time\_wind\_factor} & $k_w$ & -- & 非负 & 0.10 & 峰值风速修复放大\\
\fld{max\_repair\_time\_hr} & $T_{max}$ & h & 正数 & 24.0 & 修复时长上限\\
\fld{staged\_post\_disaster\_repair} & -- & -- & bool & true & 启用分阶段修复\\
\fld{post\_disaster\_repair\_delay\_hr} & -- & h & 非负 & 1.0 & 最后故障后准备延迟\\
\fld{repair\_crew\_time\_per\_branch\_hr} & -- & h & 非负 & 1.0 & 每支路维修槽\\
\fld{apply\_pv\_wind\_derating} & -- & -- & bool & true & 应用 PV 风致降额\\
\fld{pv\_cutout\_start\_ms} & -- & m/s & 非负 & 25.0 & PV 降额起点\\
\fld{pv\_cutout\_full\_ms} & -- & m/s & 大于起点 & 50.0 & PV 最大降额点\\
\fld{pv\_max\_reduction} & -- & pu & $0\sim1$ & 0.40 & PV 最大降额\\
\fld{pv\_transition\_hours} & -- & h & $\ge0$ & 3 & 时间平滑\\
\end{paramtable}

\subsection{TyphoonSegmentFragility}
\compmeta{TyphoonSegmentFragility}{include/hacdcpf/analysis/typhoon\_resilience.hpp}
\begin{paramtable}
\fld{kind} & -- & -- & enum & \shortstack[l]{Overhead\\Line} & 架空/电缆/混合\\
\fld{design\_wind\_ms} & $V_d$ & m/s & 正数 & 30.0 & 风暴露阈值\\
\fld{collapse\_wind\_ms} & $V_c$ & m/s & 正数 & 55.0 & severity 饱和阈值\\
\fld{lognormal\_median\_wind\_ms} & $V_{50}$ & m/s & 正数 & 45.0 & 易损中位风速\\
\fld{lognormal\_sigma} & $\sigma$ & -- & 正数 & 0.25 & 对数标准差\\
\fld{design\_rain\_mm\_hr} & $R_d$ & mm/h & 正数 & 60.0 & 风雨暴露雨阈值\\
\fld{seg\_a\_wind} & $a$ & -- & 实数 & 3.0 & 风暴露系数\\
\fld{seg\_b\_rain} & $b$ & -- & 实数 & 1.0 & 雨暴露系数\\
\fld{seg\_c\_bias} & $c$ & -- & 实数 & -7.0 & 暴露偏置\\
\fld{curve\_number\_cn} & $CN$ & -- & $1\sim99$ & 75.0 & 电缆径流曲线数\\
\fld{manning\_n} & $n$ & -- & 正数 & 0.035 & Manning 系数\\
\fld{slope\_s0} & $S_0$ & -- & 正数 & 0.01 & 坡度\\
\fld{soil\_depth\_z\_m} & $Z$ & m & 正数 & 1.0 & 土层深度\\
\fld{porosity\_phi} & $\phi$ & pu & 正数 & 0.45 & 孔隙率\\
\fld{initial\_moisture\_theta} & $\theta_0$ & pu & $0\sim\phi$ & 0.25 & 初始含水量\\
\fld{base\_cable\_failure\_p} & $p_{c,0}$ & pu & $0\sim1$ & $10^{-4}$ & 电缆基础失效概率\\
\end{paramtable}

\subsection{TyphoonCatalogOptions}
\compmeta{TyphoonCatalogOptions}{include/hacdcpf/analysis/typhoon\_resilience.hpp}
\begin{paramtable}
\fld{samples\_per\_month} & -- & 个/月 & 正整数 & 100 & 目录采样预算\\
\fld{first\_month} & -- & 月 & $1\sim12$ & 1 & 起始月\\
\fld{last\_month} & -- & 月 & $1\sim12$ & 12 & 结束月\\
\fld{base\_seed} & -- & -- & uint & 203000 & 派生样本种子\\
\fld{horizon\_hours} & -- & h & 正数 & 48 & 每样本时域\\
\fld{time\_step\_hr} & -- & h & 正数 & 1.0 & 每样本步长\\
\fld{stochastic} & -- & -- & bool & true & 随机轨迹\\
\fld{use\_month\_defaults} & -- & -- & bool & true & 月度默认\\
\fld{use\_sst\_resource} & -- & -- & bool & true & SST 资源\\
\fld{sst\_resource\_path} & -- & -- & 路径 & 空 & SST 文件\\
\fld{catalog\_path} & -- & -- & 路径 & 空 & 目录持久化路径\\
\end{paramtable}

\subsection{TyphoonTrafficImpactOptions}
\compmeta{TyphoonTrafficImpactOptions}{include/hacdcpf/analysis/typhoon\_traffic\_impact.hpp}
\begin{paramtable}
\fld{num\_steps} & $T$ & 步 & 正整数 & 24 & 交通影响步数\\
\fld{time\_step\_hr} & $\Delta t$ & h & 正数 & 0.25 & 步长\\
\fld{start\_hour} & -- & h & 实数 & 0 & 起始轨迹小时\\
\fld{coordinate\_mode} & -- & -- & enum & Auto & Auto/Geographic/Affine\\
\fld{study\_center\_latitude} & -- & deg & 合法纬度 & 22.75 & affine 中心\\
\fld{study\_center\_longitude} & -- & deg & 合法经度 & 113.55 & affine 中心\\
\fld{affine\_network\_span\_km} & -- & km & 正数 & 20.0 & affine 空间跨度\\
\fld{auto\_geographic\_track\_distance\_km} & -- & km & 正数 & 500.0 & Auto 地理判定阈值\\
\fld{rainfall\_runoff\_coefficient} & $c_r$ & -- & 非负 & 0.85 & 全局雨径流系数\\
\fld{drainage\_rate\_mm\_hr} & $d$ & mm/h & 非负 & 25.0 & 全局排水率\\
\fld{maximum\_surface\_water\_mm} & $W_{max}$ & mm & 正数 & 500.0 & 水深上限\\
\fld{drainage\_rate\_by\_link\_mm\_hr} & -- & mm/h & map & 空 & 逐 link 排水覆盖\\
\fld{runoff\_multiplier\_by\_link} & -- & -- & map & 空 & 逐 link 径流覆盖\\
\fld{depth\_speed\_quadratic} & $a$ & -- & 实数 & 0.0009 & 深度速度二次项\\
\fld{depth\_speed\_linear} & $b$ & -- & 实数 & -0.5529 & 一次项\\
\fld{depth\_speed\_intercept\_km\_hr} & $c$ & km/h & 正数 & 86.9448 & 零水深速度\\
\fld{minimum\_open\_speed\_factor} & $f_{v,min}$ & pu & $(0,1]$ & 0.05 & 开放 link 速度下限\\
\fld{flood\_closure\_depth\_mm} & $W_c$ & mm & 非负 & 300.0 & 积水闭锁阈值\\
\fld{wind\_capacity\_reduction\_start\_ms} & $V_s$ & m/s & 非负 & 15.0 & 风容量降额起点\\
\fld{wind\_capacity\_reduction\_full\_ms} & $V_f$ & m/s & 大于起点 & 35.0 & 完全降额点\\
\fld{minimum\_wind\_capacity\_factor} & $f_{w,min}$ & pu & $[0,1]$ & 0.65 & 风容量下限\\
\fld{close\_on\_extreme\_wind} & -- & -- & bool & false & 风闭锁开关\\
\fld{wind\_closure\_ms} & $V_c$ & m/s & 非负 & 45.0 & 风闭锁阈值\\
\fld{apply\_free\_flow\_time\_profile} & -- & -- & bool & true & 写回时耗 profile\\
\fld{apply\_capacity\_profile} & -- & -- & bool & true & 写回容量 profile\\
\fld{apply\_availability\_profile} & -- & -- & bool & true & 写回可用性 profile\\
\fld{capacity\_speed\_exponent} & $\gamma$ & -- & 非负 & 1.0 & 速度到容量指数\\
\fld{minimum\_open\_capacity\_factor} & $f_{c,min}$ & pu & $(0,1]$ & 0.02 & 开放 link 容量下限\\
\end{paramtable}

\subsection{核心结果结构}
\begin{center}\footnotesize
\begin{longtable}{@{}L{4.8cm}L{4.4cm}L{5.0cm}@{}}
\toprule
结构 & 字段 & 语义\\
\midrule\endfirsthead
\toprule 结构 & 字段 & 语义\\\midrule\endhead
\fld{ScenarioCandidate} & id, family, probability & 稳定候选 ID、族、组内概率\\
 & time\_series & 实际 MW/MWh/SOC profiles\\
 & contingency, resilience\_event & 条件元数据；按族可选\\
 & features, outage\_signature, regime\_label & 聚类输入；签名为内部统一目录空间\\
 & risk\_score, risk\_source\_scores & 综合及分源风险摘要\\
 & anchor\_reasons, tail\_anchor, frozen\_medoid & 风险原因与优化状态\\
 & component\_load\_profiles & 稳定组件 ID 归因；大案例可 aggregate fallback\\
 & standard\_time\_series & 下游 multiplier 契约\\
\midrule
\fld{ScenarioCluster} & cluster\_id, representative\_id & 簇号与代表候选 ID\\
 & representative & 完整代表候选\\
 & probability & 成员概率和，再归一\\
 & member\_count, member\_ids & 成员审计\\
 & average\_distance, max\_distance & 加权平均距离与半径\\
 & frozen\_medoid, anchor\_reasons & 锚点状态\\
\midrule
\fld{RegularScenarioResult} & \fld{candidate\_count}, \fld{cluster\_count}, clusters, audit & 所有 SSP/year 组合总计\\
\fld{ReliabilityScenarioResult} & \fld{contingency\_count}, \fld{cluster\_total}, contingencies, audit & 条件组嵌套结果\\
\fld{ResilienceScenarioResult} & \fld{intensity\_count}, \fld{cluster\_total}, intensities, audit & 强度组嵌套结果\\
\fld{ScenarioGenerationResult} & regular, reliability, resilience, warnings, summary & 组合入口结果\\
\midrule
\fld{TyphoonFaultSequenceResult} & track, generated\_segments, faults & 轨迹/分段/故障\\
 & \fld{branch\_indices}, branch peak arrays, \fld{branch\_risks} & 稳定支路 ID 与峰值风险\\
 & \fld{renewable\_profile\_multiplier} & 灾害新能源倍率\\
 & used\_* flags & 坐标、合成段、目录、类别、修复近似\\
 & requested/selected category, sample metadata & 请求与实际目录样本\\
 & \fld{monthly\_sst\_c}, \fld{sst\_resource\_path}, seed, status & 资源与状态\\
\midrule
\fld{TyphoonTrafficImpactResult} & impacted\_traffic, link\_steps & 写回图与逐 link/step 结果\\
 & \fld{links\_closed\_at\_least\_once}, \fld{closed\_link\_steps} & 闭锁摘要\\
 & \fld{used\_affine\_georeferencing} & 坐标 fallback\\
 & peak wind/rain/water, minimum factors & 聚合极值\\
 & model\_scope & 模型覆盖声明\\
\bottomrule
\end{longtable}
\end{center}

\subsection{条件、线段、风险与逐步结果逐字段表}
以下字段是前表中“条件元数据”和“逐步结果”的完整展开，用于保证公开头字段可以按名称检索。
\begin{center}\footnotesize
\begin{longtable}{@{}L{4.5cm}L{4.7cm}L{5.0cm}@{}}
\toprule
结构 & 字段 & 语义/单位\\
\midrule\endfirsthead
\toprule 结构 & 字段 & 语义/单位\\\midrule\endhead
\fld{ContingencyDefinition} & \fld{display\_name} & 人类可读名称，不作为身份\\
 & \fld{affected\_ac\_branches} & 带域前缀的稳定 AC 支路引用\\
 & \fld{affected\_dc\_branches} & 带域前缀的稳定 DC 支路引用\\
 & \fld{affected\_converters} & 受影响 VSC 等换流器引用\\
 & \fld{affected\_dcdc\_converters} & 受影响 DC/DC 引用\\
 & \fld{affected\_buses} & 受影响母线引用，必须保留域\\
 & \fld{affected\_generators} & 受影响发电资产引用\\
 & \fld{affected\_loads} & 受影响负荷引用\\
 & \fld{affected\_storage} & 受影响储能引用\\
\midrule
\fld{ResilienceEventDefinition} & \fld{repair\_fault\_count} & 修复排序处理的故障数\\
 & \fld{pv\_typhoon\_multiplier} & 逐步 PV 台风倍率\\
 & \fld{wind\_typhoon\_multiplier} & 逐步风电台风倍率\\
 & \fld{renewable\_typhoon\_multiplier} & 容量加权新能源倍率\\
 & \fld{load\_typhoon\_multiplier} & 逐步负荷台风倍率\\
 & \fld{stage1\_start\_step} & 可选第一阶段起始步\\
 & \fld{stage1\_end\_step} & 可选第一阶段终止步\\
\midrule
\fld{TyphoonLineSegment} & \fld{branch\_kind} & AC/DC 支路域\\
 & \fld{branch\_index} & 稳定支路 index\\
 & \fld{segment\_index} & 支路内 0-based 段号\\
 & \fld{from\_latitude}, \fld{from\_longitude} & 段起点坐标（deg）\\
 & \fld{to\_latitude}, \fld{to\_longitude} & 段终点坐标（deg）\\
 & \fld{mid\_latitude}, \fld{mid\_longitude} & 风雨求值中点（deg）\\
 & \fld{length\_km} & 段长（km）\\
 & \fld{fragility} & 该段完整易损参数\\
\midrule
\fld{TyphoonBranchRisk} & \fld{branch\_kind}, \fld{branch\_index} & 稳定带域支路身份\\
 & \fld{peak\_wind\_ms} & 峰值站点风速（m/s）\\
 & \fld{peak\_rain\_mm\_hr} & 峰值雨强（mm/h）\\
 & \fld{peak\_failure\_probability} & 峰值条件失效概率（pu）\\
\midrule
\fld{TyphoonCatalog} & \fld{category\_counts} & 各类别目录样本数，非发生概率\\
\midrule
\fld{TyphoonFaultSequenceResult} & \fld{branch\_peak\_wind\_ms} & 与 \fld{branch\_indices} 同序，m/s\\
 & \fld{branch\_peak\_rain\_mm\_hr} & 与 \fld{branch\_indices} 同序，mm/h\\
 & \fld{branch\_peak\_failure\_probability} & 与 \fld{branch\_indices} 同序，pu\\
 & \fld{repair\_fault\_count} & 参与修复排序的故障数\\
\midrule
\fld{TyphoonTrafficLinkStep} & \fld{midpoint\_latitude}, \fld{midpoint\_longitude} & link 暴露求值中点（deg）\\
 & \fld{free\_flow\_speed\_factor} & 深度速度因子（pu）\\
 & \fld{available} & 该 link-step 是否开放\\
\midrule
\fld{TyphoonTrafficImpactResult} & \fld{peak\_wind\_ms} & 全 link-step 峰值风速（m/s）\\
 & \fld{peak\_rainfall\_mm\_hr} & 全 link-step 峰值雨强（mm/h）\\
 & \fld{peak\_surface\_water\_mm} & 峰值积水深度（mm）\\
 & \fld{minimum\_speed\_factor} & 最小开放速度因子（pu）\\
 & \fld{minimum\_capacity\_factor} & 最小容量因子（pu）\\
\midrule
\fld{ScenarioBundleIoReport} & \fld{sheet\_counts} & 工作簿表名与记录数\\
\bottomrule
\end{longtable}
\end{center}

\subsection{结果解释纪律}
\begin{itemize}
  \item \fld{success=true} 只表示序列化了生成结果，不表示下游 PF/OPF 可行；
  \item \fld{candidate\_count} 与 \fld{cluster\_count} 是实际计数；可能因多组合而大于单配置；
  \item \fld{audit} 是模型证据的一部分，不能在 HTTP 层丢弃；
  \item \fld{warnings}、fallback/approximate 标志必须随导出保留；
  \item 结果向量的支路/组件引用必须通过稳定 ID 解释，不能按 vector position 对号入座。
\end{itemize}

\section{JSON、HTTP 与工作簿契约}
\label{sec:scenario-io-contract}

\subsection{选项 JSON}
\srcpath{scenario\_generation\_options\_from\_json} 读取顶层
\texttt{regular}、\texttt{reliability}、\texttt{resilience}、\texttt{perturbation}、
\texttt{typhoon\_impact} 和 \texttt{clustering}。缺失字段使用结构默认值；相关系数、分位数、
边界宽度和非负权重按第~\ref{sec:scenario-options-results} 章钳制。可靠性
\fld{num\_steps} 最终强制为 1。

最小请求示例：
\begin{verbatim}
{
  "regular": {"enabled": true, "candidate_count": 128,
              "cluster_count": 12, "ssp": "ssp245", "year": 2050},
  "reliability": {"enabled": false},
  "resilience": {"enabled": false},
  "perturbation": {"seed": 20260822, "temporal_correlation": 0.75,
                   "load_renewable_correlation": -0.25},
  "clustering": {"method": "hybrid_kmedoids_tail_5pct",
                 "boundary_fraction": 0.10, "compare_baseline": true}
}
\end{verbatim}

\subsection{结果 JSON}
\srcpath{scenario\_generation\_result\_to\_json} 返回
\texttt{success}、\texttt{summary}、\texttt{warnings}、\texttt{regular}、
\texttt{reliability} 和 \texttt{resilience}。每族保留实际计数、nested groups、clusters 和
audit。cluster 内嵌完整 representative candidate，因此 time series、条件元数据、风险、锚点、
成员和距离可以沿一条 JSON 路径闭环，不需依赖 GUI 内存。

\subsection{生产 HTTP：生成场景}
生产路由位于 \srcpath{tests/run\_gui\_server.cpp}：
\begin{verbatim}
POST /api/session/generate_scenarios
Content-Type: application/json
\end{verbatim}
路由先取得不可变 session system snapshot；已有分析运行时返回 HTTP 409。解析选项后，HTTP 层
有意强制：
\begin{equation}
 T_{regular}=8760,\qquad T_{reliability}=1,\qquad T_{resilience}=48.
\end{equation}
预算钳制为 regular candidates $[1,500]$、reliability/resilience candidates $[1,200]$，
各 cluster count 在 1 到对应候选数之间。类别覆盖 cluster count 同样钳制。故而直接 C++ API
允许的短时域 review case 与 HTTP 生产路由不是相同行为。

响应在公共 JSON 外增加 \texttt{\_performance}：generation ms 和三族 load processing
granularity；HTTP header 还返回 \texttt{Server-Timing} 与
\texttt{X-Scenario-Response-Bytes}。这些是性能元数据，不是算法收敛证书。

\subsection{生产 HTTP：台风故障预览}
\begin{verbatim}
POST /api/session/generate_typhoon_faults
\end{verbatim}
该路由读取 horizon、step、seed、stochastic、month、SST path、segment length、sampling gate、
staged repair 等公开子集，生成线路风险和故障供前端预览。未在该路由读取的
\fld{TyphoonScenarioOptions} 字段只能通过 C++ API 或后续接口使用；不能因结构公开就声称 HTTP
已全量暴露。

\subsection{场景包 v3}
规范包外壳为
\begin{verbatim}
{
 "format":"generated_scenario_case_bundle_v3",
 "schema_version":3,
 "unit_space":"dimensionless_multiplier",
 "family":"mixed",
 "case_count":1,
 "cases":[{
   "case_id":"regular:ssp245:2050:17",
   "system":{...},
   "generated_scenario":{...},
   "standard_time_series":{"num_steps":8760,"profiles":[...]}
 }]
}
\end{verbatim}
v3 的 profile 值是 multiplier 或 soc fraction；设备表另存 base MW/MWh 和 capacity field，
避免把倍率当有名功率。\texttt{system} 中资产仍使用有名值。

\subsection{规范化}
\srcpath{normalize\_generated\_scenario\_bundle} 支持三种输入：
\begin{enumerate}
  \item 已是 v3：原样返回；
  \item 单个 system，含 \texttt{\_generated\_scenario}/\texttt{\_time\_series}：包装为一个 case；
  \item v1/v2 cases 数组：规范化 case shape，升级 format/schema/unit space，并产生 warning。
\end{enumerate}
不含 system 且不含 cases 的对象抛出异常。升级 warning 是数据迁移证据，调用方不应丢弃。

\subsection{Strict 与 Permissive 验证}
\srcpath{validate\_scenario\_bundle} 检查：
\begin{itemize}
  \item \texttt{unit\_space=dimensionless\_multiplier}；cases 非空；
  \item profile ID 不重复，values 长度等于 num steps；
  \item 所有值有限；SOC 在 $[0,1]$；其他 multiplier 非负；
  \item case ID 和标准时序结构可解析。
\end{itemize}
Strict 把问题写入 \fld{errors}，Permissive 写入 \fld{warnings}；两者均返回规范化内容，调用方
必须检查 report，不能仅凭函数返回判为有效。

验证路由：
\begin{verbatim}
POST /api/session/validate_scenario_bundle
\end{verbatim}
生产路由当前使用 Permissive，响应显式返回 bundle、warnings、errors 与 case count/unit space。

\subsection{工作簿布局}
启用 \texttt{HACDCPF\_ENABLE\_ETAP} 时，OpenXLSX 路径写入 Metadata、Cases、Devices、Profiles
等表。Metadata 保存 workbook schema、JSON format/schema、unit space、family、case count；
Devices 保存 case ID、kind、稳定 component index、position（仅校验）、bus、name、base MW/MWh、
capacity field、service 和 profile ID；Profiles 保存逐时 multiplier。

导入优先按 kind+稳定 index 定位设备，找不到才使用 position fallback。position 回退必须形成
warning，不能成为跨版本稳定身份。

\subsection{工作簿 HTTP 与编译门}
\begin{verbatim}
POST /api/session/export_scenario_workbook
POST /api/session/import_scenario_workbook
\end{verbatim}
导出接收 JSON bundle、返回 xlsx binary；导入接收 binary、返回 bundle/report。临时文件在异常
路径也清理。未编译 \texttt{HACDCPF\_ENABLE\_ETAP} 时，公共 workbook API 明确抛出
\texttt{Scenario workbook I/O not compiled}，不是返回空文件。

\subsection{合同边界}
GUI 下载成功只证明 HTTP 和文件 I/O 成功，不证明 scenario probability、气候资源或 PF/OPF
可行。v3 校验只验证结构、有限数和 multiplier 语义，不重新运行候选生成或聚类。

\section{验证架构、测试目录与外部对照边界}
\label{sec:scenario-validation-architecture}

\subsection{验证金字塔}
\begin{center}\small
\begin{tabular}{@{}L{3.2cm}L{5.3cm}L{5.8cm}@{}}
\toprule
层级 & 证据 & 能证明的对象\\
\midrule
L1 解析/不变量 & 独立公式、概率和、有限数 & 局部方程与守恒\\
L2 单元回归 & Catch2 固定 seed/边界案例 & 分支、fallback、字段消费\\
L3 交叉实现 & C++ driver + 独立 Python & 序列化结果可由公开方程复算\\
L4 接口/模式 & JSON v3、workbook、HTTP & 数据契约和生产路由\\
L5 下游任务 & PF/OPF/reliability/resilience 回测 & 代表场景对具体决策任务的充分性\\
\bottomrule
\end{tabular}
\end{center}
当前模块具备 L1--L4 的专项证据；L5 需要按下游模块另行设计，不能由场景生成测试代替。

\subsection{test\_scenario\_generation 测试目录}
\begin{longtable}{@{}L{8.0cm}L{6.0cm}@{}}
\toprule
Catch2 case & 验证对象\\
\midrule\endfirsthead
\toprule Catch2 case & 验证对象\\\midrule\endhead
Regular scenario generation applies SSP/year climate morphing to load and PV & 已知 SSP/year 气候形变\\
Scenario generation prepares standard multiplier time series internally & multiplier/base/capacity 契约\\
Regular climate perturbation creates reproducible candidate variation & 固定 seed 与气候随机性\\
Scenario perturbations realize declared temporal and cross correlations & AR(1) 与跨变量统计目标\\
Hybrid clustering consumes the declared quantile boundary fraction & boundary fraction 字段消费\\
Typhoon minimum fault probability gates sampling without hiding risk & sampling gate 与风险保留\\
Regular climate morphing falls back safely for unknown SSP/year & 零增量 fallback/audit\\
Large annual scenario generation uses bounded aggregate load profiles & 年度内存预算与 aggregate 降级\\
Large reliability catalog bounds audit samples without dropping contingencies & audit 预算不删目录\\
Very large reliability catalogs preserve N-1 coverage within work budgets & 大目录 N--1 覆盖\\
Typhoon fault generation reuses precomputed topology segments & 空间预计算复用\\
Large typhoon fault sets use bounded upstream-first repair ordering & 近似修复排序与标志\\
\bottomrule
\end{longtable}

\subsection{test\_typhoon\_traffic\_impact 测试目录}
\begin{longtable}{@{}L{8.0cm}L{6.0cm}@{}}
\toprule
Catch2 case & 验证对象\\
\midrule\endfirsthead
Typhoon road impact creates common speed, capacity, and closure profiles & 逐 link profile 与闭锁\\
Typhoon road impact marks affine georeferencing of benchmark topology & Auto/affine fallback 标志\\
CTM reads road profiles by parent simulation step & 下游时步映射\\
Typhoon catalog snapshots survive concurrent mixed-option refreshes & 缓存线程安全快照\\
Typhoon road impact rejects invalid clamp and hydrology bounds & 选项失败语义\\
\bottomrule
\end{longtable}

\subsection{test\_scenario\_bundle\_schema 测试目录}
\begin{longtable}{@{}L{8.0cm}L{6.0cm}@{}}
\toprule
Catch2 case & 验证对象\\
\midrule\endfirsthead
Scenario bundle validation accepts v3 dimensionless multiplier schema & v3 正常路径\\
Scenario bundle validation rejects invalid multiplier profiles & 负值/SOC/长度等错误\\
Scenario bundle normalization upgrades v2 case shape & v2→v3 迁移与 warning\\
\bottomrule
\end{longtable}

\subsection{独立 oracle 设计}
\srcpath{tests/scenario\_generation\_cross\_validation.py} 不链接也不导入 hacdcpf。它从 C++ JSON
读取公开输入，独立实现 Haversine、Holland、降雨、等效风、架空段失效、交通积水/容量、
Pearson 统计、median/IQR 缩放和簇距离复算。验证项包括：
\begin{itemize}
  \item 候选和代表概率误差 $<10^{-12}$；
  \item 4096 步 load lag-1 与目标差 $<0.05$；全天 load-wind 同期相关与目标差 $<0.08$；
  \item 风、雨、易损概率和交通递推误差 $<10^{-10}$；
  \item 128 个成员无遗漏、代表属于成员、平均/最大距离和 transport 指标误差 $<10^{-10}$；
  \item sampling gate 下 fault count 为 0 且风险概率仍保留。
\end{itemize}

\subsection{为什么 cross-correlation 使用风电}
PV 基线夜间为零，只能在白天恢复 multiplier。第一次预声明容差为 0.08 时，PV 子样本实测
-0.341279，相对目标 -0.25 偏差 0.091279，验证失败。没有放宽阈值，而是重新推导实验：使用
全天非零 wind profile 恢复完整 renewable multiplier，实测 -0.290252，偏差 0.040252，通过
原阈值；PV 值继续报告为周期性子样本诊断。这是测试设计修复，不是结果筛选。

\subsection{OpenDSS 与 GridLAB-D 的正确角色}
OpenDSS 和 GridLAB-D 能消费确定性负荷/发电 profile 并运行时序潮流，但它们没有与本模块
等价的 SSP 气候形变、台风目录、易损故障抽样和尾部锚定 k-medoids API。因此：
\begin{itemize}
  \item 可以把代表 profile 导出后比较下游电压/潮流；
  \item 不能把下游潮流一致写成“场景生成算法认证”；
  \item 本手册没有伪造两软件的聚类或灾害链数值；
  \item 若未来增加下游 profile replay，应单独记录软件版本、网络转换、控制模式和误差。
\end{itemize}

\subsection{统计准入与防止移动门槛}
容差必须在运行前声明。若实测偏差超限，按以下顺序处理：
\begin{enumerate}
  \item 检查 oracle 是否独立、输入是否相同、单位/索引是否一致；
  \item 检查有限样本、AR 有效样本量、clip 和周期性选择；
  \item 检查 C++ 是否忠实实现理论；
  \item 只有理论重新推导支持时才修改实验或门槛，并保留失败原因。
\end{enumerate}

\subsection{未覆盖测试}
当前尚无：历史台风事件的空间观测评分、校准易损曲线的 out-of-sample 检验、场景数序贯收敛、
下游 OPF 决策 regret、跨平台 bitwise 重现、v3 工作簿大规模 round-trip 性能基线。这些必须
列为未认证边界，不能写入“已全面验证”。

\section{固定数值案例、统计结果与独立交叉验证}
\label{sec:scenario-numerical-validation}

\subsection{案例系统与复现条件}
\srcpath{tools/scenario\_generation\_review\_case.cpp} 构造 2 母线 AC 系统：slack bus 1、load bus 2、
支路 \texttt{ac\_branch:101}、10 MW/2 Mvar 负荷 201、5 MW PV 301、3 MW wind 302 和
20 MW 上限发电机 401。母线坐标为 $(22.75,113.55)$ 与 $(22.75,113.95)$。

常规短时案例参数：24 步、64 候选、6 代表、seed 20260822、load sigma 0.08、renewable sigma
0.12、$\rho_t=0.75$、$\rho_{LR}=-0.25$、block=1、无气候/SOC 扰动。未知
\texttt{ssp999:2090} 有意触发零气候增量 fallback。统计长序列用 4096 步、1 候选，sigma 改为
0.05/0.08，保持同一相关参数。

缩减案例使用 128 候选、12 代表、hybrid tail-anchor、seed 17、
\fld{boundary\_fraction}=0.10、\fld{regime\_weight}=0、baseline comparison=true。

\subsection{概率与确定性复现}
\begin{center}\small
\begin{tabular}{@{}L{6.4cm}L{4.2cm}L{3.5cm}@{}}
\toprule
量 & 数值 & 判据\\
\midrule
常规候选数 / 代表数 & 64 / 6 & 等于请求\\
候选概率和 & 1.000000000000 & 误差 0\\
代表概率和 & 1.000000000000 & 误差 0\\
重复生成 coverage/簇数 & 完全相同 & deterministic=true\\
64 候选 load--PV energy correlation & -0.3663644365 & 目标 -0.25，容差 0.15\\
\bottomrule
\end{tabular}
\end{center}
能量相关同时含确定性日内 profile、PV 夜间零值、风电容量和有限候选，因此只作为端到端统计
现象；AR 核的直接验证使用下面的 multiplier 分离实验。

\subsection{4096 步 AR(1) 结果}
用同一系统和同一确定性生成器分别运行“有扰动”和“无扰动”，逐点相除恢复 multiplier。
负荷和风电基线始终为正；PV 只保留 1879 个白天点。
\begin{center}\small
\begin{tabular}{@{}L{5.7cm}rrrr@{}}
\toprule
序列 & count & mean & sample std & min / max\\
\midrule
load multiplier & 4096 & 0.99804200 & 0.04915002 & 0.83864293 / 1.17151450\\
wind multiplier & 4096 & 1.00088103 & 0.08094124 & 0.73517822 / 1.27080447\\
PV daylight multiplier & 1879 & 1.00471410 & 0.08273485 & 0.73517822 / 1.27080447\\
\bottomrule
\end{tabular}
\end{center}

相关结果：
\begin{center}\small
\begin{tabular}{@{}L{6.8cm}rrr@{}}
\toprule
统计量 & 目标 & 实测 & 绝对偏差\\
\midrule
load lag-1 correlation & 0.750000 & 0.74375318 & 0.00624682\\
load--wind contemporaneous & -0.250000 & -0.29025168 & 0.04025168\\
load--PV daylight contemporaneous & -0.250000 & -0.34127948 & 0.09127948\\
\bottomrule
\end{tabular}
\end{center}
预声明门为 lag-1 0.05、同期 0.08。全天 wind 通过；PV 白天子样本不作为准入门并保留失败数值。
这避免通过事后放宽容差掩盖周期性筛选偏差。

\subsection{128→12 场景缩减结果}
\begin{center}\small
\begin{tabular}{@{}L{5.5cm}rrr@{}}
\toprule
指标 & Hybrid & Weighted baseline & Hybrid--baseline\\
\midrule
transport mean & 2.14810121 & 1.18838011 & +0.95972110\\
transport max & 7.16929262 & 3.06119436 & +4.10809826\\
global tail coverage & 1.000000 & 0.000000 & +1.000000\\
coupling tail coverage & 0.450000 & 0.150000 & +0.300000\\
load source coverage & 0.714286 & 0.142857 & +0.571429\\
net-load source coverage & 0.571429 & 0.285714 & +0.285714\\
PV deficit coverage & 0.714286 & 0.285714 & +0.428571\\
renewable deficit coverage & 0.857143 & 0.285714 & +0.571429\\
wind deficit coverage & 1.000000 & 0.000000 & +1.000000\\
regime mass L1 & 1.640625 & 0.046875 & +1.593750\\
\bottomrule
\end{tabular}
\end{center}
Hybrid 的 12 个 anchor 全部冻结。结论是明确的双目标取舍：尾部覆盖大幅提高，但平均/最坏运输
误差和 regime 质量显著变差，不能写成“hybrid 全面优于 baseline”。

\subsection{聚类独立复算}
Python oracle 从 C++ 序列化的 128 个候选和 12 个簇重新：
\begin{enumerate}
  \item 按字段列构造 median/IQR 尺度；
  \item 检查每个 representative ID 属于 member IDs；
  \item 逐成员计算欧氏距离（案例 regime weight=0，无 outage 差异）；
  \item 复算 cluster probability、weighted average distance、max distance；
  \item 汇总 transport mean/max 和成员总数。
\end{enumerate}
结果：成员总数 128、概率误差 0、最大平均距离误差 0、最大半径误差 0、transport mean/max 误差
均为 0（断言容差 $10^{-10}$）；tail/frozen medoid count 均为 12。

\subsection{固定 Holland、雨量与易损结果}
固定风暴参数：中心 $(22.75,113.55)$，$\Delta p=70$ hPa，$B=1.2$，$R_{mw}=40$ km，
$f_c=5.6398700067\times10^{-5}$ s$^{-1}$，站点 $(22.75,113.95)$。结果：
\begin{center}\small
\begin{tabular}{@{}L{6.4cm}L{4.8cm}L{3.0cm}@{}}
\toprule
量 & C++ 值 & Python 误差\\
\midrule
Holland 10 m wind & 55.24244506292464 m/s & 0\\
rainfall & 32.58499937279819 mm/h & $7.82\times10^{-14}$\\
2 km overhead failure probability & 0.7989101008439494 & $6.66\times10^{-16}$\\
\bottomrule
\end{tabular}
\end{center}
sampling gate 设为 1.0，所以 fault count=0；风险 0.798910 仍在 branch risks，验证门不隐藏风险。

\subsection{交通递推结果}
固定 link 长 4 km、free-flow time 0.08 h、capacity 1200，runoff=0.85、drainage=25 mm/h、
$\Delta t=1$ h、闭锁水深 300 mm。逐步结果：
\begin{center}\small
\begin{tabular}{@{}rrrrrr@{}}
\toprule
step & rain & water & speed factor & wind factor & capacity factor\\
\midrule
0 & 32.584999 & 2.697249 & 0.982923 & 0.65 & 0.638900\\
1 & 32.584999 & 5.394499 & 0.965996 & 0.65 & 0.627898\\
2 & 32.584999 & 8.091748 & 0.949221 & 0.65 & 0.616993\\
3 & 32.584999 & 10.788998 & 0.932595 & 0.65 & 0.606187\\
\bottomrule
\end{tabular}
\end{center}
closed link-step=0；独立水深最大误差 $6.39\times10^{-14}$。

\subsection{可靠性目录结果}
review system 的目录为 \texttt{ac\_gen:401}、\texttt{ac\_branch:101}、
\texttt{ac\_ren:302}，count=3，count 与 ID 数组长度一致。PVSystem 当前没有由本案例 include
组合产生独立目录项；文档按实际输出报告，不用 enum 覆盖列表替代数值结果。

\subsection{证据文件与复现命令}
证据目录为 \srcpath{external\_data/scenario\_generation\_validation/}。原始 C++ 输出文件为
\srcpath{scenario\_generation\_review\_case.json}；独立报告文件为
\srcpath{scenario\_generation\_cross\_validation.json}。
\begin{verbatim}
cmake --build build/macos-release --target scenario_generation_review_case -j4
python3 tests/scenario_generation_cross_validation.py \
  --cpp-bin build/macos-release/tests/scenario_generation_review_case \
  --output-dir external_data/scenario_generation_validation
\end{verbatim}

\subsection{结论边界}
这些数值证明公开方程、字段消费、概率守恒、聚类序列化和固定 seed 复现一致；它们不证明气候
或台风预测准确、易损参数经现场校准、下游 PF 可行或代表场景在任意 OPF 目标下最优。

\section{系统审计、准入规则与未闭合边界}
\label{sec:scenario-audit-admission}

\subsection{理论--实现--结果闭环矩阵}
\begin{longtable}{@{}L{3.6cm}L{4.0cm}L{3.6cm}L{3.2cm}@{}}
\toprule
对象 & 实现锚点 & 数值/测试证据 & 状态\\
\midrule\endfirsthead
\toprule 对象 & 实现锚点 & 数值/测试证据 & 状态\\\midrule\endhead
族内条件概率 & 三族 generate 函数 & probability error=0 & 已闭环\\
二元 AR(1) & correlated noise factors & 4096 步 lag/cross & 已闭环\\
block interpolation & 同上 & 单元统计测试 & 已闭环；逐时参数有偏差\\
气候形变 & climate morph 函数族 & known/unknown SSP tests & 已闭环；非气候预报\\
component profiles & candidate builder & standard TS test & 已闭环；大案例 aggregate fallback\\
储能 SOC & storage SOC profiles & profile/schema tests & 已闭环；时域内常值\\
N--1 目录 & enumerate contingencies & 大目录与 review IDs & 已闭环；无故障率先验\\
台风目录 & catalog build/sample & 并发 snapshot test & 已闭环；无历史概率校准\\
Holland/雨量 & wind/rain kernels & Python error $<10^{-10}$ & 已闭环\\
架空易损 & overhead probability & Python error $6.7\times10^{-16}$ & 已闭环；默认参数未校准\\
电缆水文 & cable probability & 单元分支覆盖 & 公式闭环；缺独立固定点表\\
sampling gate & fault sequence & risk retained, fault=0 & 已闭环\\
staged repair & repair ordering & 大故障集 test & 已闭环；非全局最优\\
交通积水/容量 & traffic impact & Python error $<10^{-10}$ & 已闭环\\
tail anchors & risk metadata/select & boundary test & 已闭环\\
k-medoids & clustering core & 128→12 独立复算 & 已闭环；局部最优\\
JSON v3 & bundle normalize/validate & 3 schema tests & 已闭环\\
HTTP & production routes & 路由源码合同 & 已实现；缺专项 E2E 数值\\
workbook & scenario bundle I/O & schema/编译门 & 已实现；大规模 round-trip 待测\\
OpenDSS/GridLAB-D & 无等价生成内核 & 未伪造数值 & 不适用直接认证\\
下游 decision regret & 无 & 无 & 未实现\\
\bottomrule
\end{longtable}

\subsection{结果准入硬门}
发布、导出或引用结果前必须同时满足：
\begin{enumerate}
  \item 每个条件组候选/代表概率有限、非负，概率和误差 $<10^{-12}$；
  \item profile 数值有限，SOC 在 $[0,1]$，multiplier 非负，长度匹配 num steps；
  \item 稳定组件/支路 ID 可回溯，AC/DC 域未混用，member 数无遗漏；
  \item 实际方法、候选数、代表数、seed、资源来源和 audit 随结果保存；
  \item 所有 fallback/approximate 标志与 warnings 未被 HTTP/导出层删除；
  \item 固定公式 oracle 在预声明 $10^{-10}$ 数值容差内；统计检验使用预声明有限样本容差；
  \item 若用于下游 PF/OPF/可靠性/弹性结论，另有该下游任务的可行性和外样本验证。
\end{enumerate}

\subsection{必须降级口径的情形}
出现以下任一情形，结果仍可用于探索，但必须降低 model scope：未知 SSP/year 零形变；缺组件
profile 使用 aggregate；缺坐标或真实线段；使用默认未校准 fragility；类别 fallback；锚点截断；
全部 medoid 冻结；修复近似排序；达到聚类迭代上限；工作簿 position fallback；下游求解不收敛。

\subsection{不得作出的声明}
\begin{itemize}
  \item 不得把固定 seed 重现称为统计收敛；
  \item 不得把族内等概率称为历史发生概率；
  \item 不得把 SSP label 称为 SSP 概率；
  \item 不得把 Holland 固定点公式吻合称为台风预报准确；
  \item 不得把默认易损曲线称为 IEC、国标或现场认证；
  \item 不得把 greedy/排序修复称为全局最优抢修；
  \item 不得把 hybrid 尾部覆盖提升称为缩减全面优越；
  \item 不得把 OpenDSS/GridLAB-D profile replay 称为场景生成内核交叉认证；
  \item 不得把 \texttt{success=true} 称为 PF/OPF 可行或工程可投产。
\end{itemize}

\subsection{已知未闭合研究边界}
本版如实保留以下边界：
\begin{enumerate}
  \item 没有多变量空间天气场、copula/VAR、极值尾部校准和概率预报评分；
  \item 没有历史台风后验、登陆率、类别/月度先验和观测风雨空间验证；
  \item 没有结构/电缆易损参数的现场数据校准与置信区间；
  \item 没有相关故障、保护/通信共因和动态修复 Markov 模型；
  \item 没有二维水动力、交通需求重分配、抢修可达性闭环；
  \item 没有 decision-aware reduction、PF/OPF regret 或自适应样本收敛；
  \item 没有对生产 HTTP 的场景生成专项大案例 E2E，也没有大工作簿性能基线。
\end{enumerate}
这些不是“已完成但未写文档”的能力；它们当前不在实现中。

\subsection{变更准入流程}
后续修改场景模块必须依次完成：
\begin{enumerate}
  \item 在理论章写出概率对象、方程、单位、假设和复杂度预测；
  \item 在实现函数附近引用方程/文献，并消费所有新增公开字段；
  \item 增加正常、边界、fallback、无效输入和确定性测试；
  \item 对可独立复算的数值核增加外部 oracle，不与主库共享实现；
  \item 更新结果 audit/model scope、JSON/HTTP/workbook 和本手册；
  \item 复跑 Catch2、CTest、静态字段审计、PDF 编译和逐页渲染；
  \item 实测偏离理论预测时执行重新推导，不移动阈值或删除负面结果。
\end{enumerate}

\subsection{本版准入结论}
本版可以声明：三族条件生成、二元相关扰动、参数化台风灾害链、交通递推、混合缩减和 v3
数据契约均有源码和测试，核心固定公式/聚类输出有独立数值复算。不能声明：气候/台风预测、
易损度现场校准、场景样本统计收敛、下游决策最优性或 OpenDSS/GridLAB-D 等价算法认证。

\section{跨模块工业级技术评价基线}

本章规定本手册所述模块进入工程算例、批量研究或生产服务前必须共同满足的评价基线。
具体模型公式、变量、约束和专用算法以正文相应章节为准；本章不扩大源码已经实现的范围。

\subsection{输入、输出、边界条件与数据结构审计}
输入审计应逐项检查有限性、单位、上下界、枚举值和引用完整性。系统对象必须先按所需层级完成
校验；AC 与 DC 母线采用域限定映射，组件稳定 \texttt{index}、向量位置和临时图索引不得混用。
输出审计必须记录单位、索引空间、模型范围、收敛或可行状态、后端、容差、迭代/节点计数和告警。
空系统、空时域、孤岛、重复 ID、零容量、退化约束与不可用可选后端均属于边界测试集。

\subsection{工程假设与数值稳定性分析}
模型使用的平衡、线性化、稳态、独立性、概率分布、时间离散或网络等值假设必须与应用场景匹配。
数值评价至少包含变量缩放、绝对/相对残差、可行性违反、条件数或枢轴质量、步长/阻尼、迭代停滞
和随机种子复现性。若采用正则化、平滑、回退、松弛或近似，应同时报告触发条件、参数值、对主指标
的敏感性以及是否改变模型口径；不得用“已返回结果”替代收敛或有效性证明。

\subsection{复杂度与资源分析}
令 $n_x$、$n_c$、$n_t$、$n_s$、$|V|$、$|E|$ 分别表示变量、约束、时段、场景、节点和边数。
报告应分别给出建模、求解、后处理的时间与内存。图遍历通常为 $O(|V|+|E|)$；逐时段/逐场景
流程至少线性增长为 $O(n_t n_s)$ 次子问题调用；稠密线性代数最坏可达 $O(n_x^3)$，稀疏方法则应
报告结构非零数、填充率和因子分解峰值内存；MILP/NLP 不给出虚假的多项式最坏界，而报告节点数、
间隙、时限和实测扩展曲线。复杂度结论必须基于固定硬件、固定后端和可复现算例族。

\subsection{异常工况处理}
非法输入应在进入求解器前返回结构化错误；不可行、无界、不收敛、时限、内存不足和后端缺失必须
相互区分。部分结果只有在对应 \fld{ValidityFlags}、\fld{model\_scope}、
\fld{model\_limitations} 或等价字段允许时才可使用。异常路径不得静默丢弃 AC/DC 资产、制造平衡
节点、把 NaN 改写为零或把近似解标为全局最优；系统原始对象应保持可恢复和可重试。

\subsection{与其他模块的接口关系}
接口链路遵循“工程语义模型 $\to$ 校验 $\to$ 投影/装配 $\to$ 求解或分析 $\to$ 结果回投影”。
跨模块传递时保留稳定 ID、单位、符号、时间基准和场景身份；消费潮流、OPF、动态或可靠性结果前，
先检查其收敛、覆盖范围和有效性标志。HTTP/JSON/Python 层只做语义保持适配，不重新解释求解结果。

\subsection{测试与验证建议}
最低验证矩阵包括：手算小例、标准算例、边界/异常输入、随机性质测试、算法或后端交叉校核、
序列化往返、稳定 ID 回投影、AC/DC 同号母线、重复运行确定性和规模扩展测试。数值算法还应加入
残差独立重算、雅可比有限差分或自动微分校核；近似模型应与高保真模型比较并给出误差分布，而非
只报告单个平均值。回归报告必须列明构建配置、算例版本、容差、通过/失败/跳过数及未验证范围。

\subsection{工业级评估指标}
\begin{itemize}
  \item \textbf{正确性}：守恒残差、约束最大违反、解析/基准误差、结果回投影一致性；
  \item \textbf{鲁棒性}：收敛率、不可行识别率、异常分类准确率、扰动敏感性与随机复现率；
  \item \textbf{性能}：端到端时延、吞吐量、峰值内存、并行效率与规模增长斜率；
  \item \textbf{可追溯性}：输入模式版本、稳定 ID 覆盖率、模型范围/限制字段完整率、告警可定位率；
  \item \textbf{工程有效性}：与标准、外部引擎或现场数据的偏差，以及误判/漏判对业务指标的影响。
\end{itemize}
指标应预先给出验收阈值和失败处置；未达到阈值时只能标记为实验性或受限适用。

\subsection{适用范围、局限性与后续改进}
本手册只评价当前源码覆盖的模型、后端和数据结构，不对未建模物理过程、未安装求解器或未执行的
外部验证作保证。后续改进应按“缺口 $\to$ 可量化指标 $\to$ 源码/测试变更 $\to$ 独立验证”闭环，
优先消除会造成错误结论的语义缺口，其次提升数值鲁棒性与性能，最后扩展模型范围。


\end{document}
